\documentclass{article}

\PassOptionsToPackage{numbers,compress}{natbib}
\usepackage{neurips_2026}

\usepackage[utf8]{inputenc}
\usepackage[T1]{fontenc}
\usepackage{hyperref}
\usepackage{url}
\usepackage{booktabs}
\usepackage{amsfonts}
\usepackage{amsmath}
\usepackage{amssymb}
\usepackage{nicefrac}
\usepackage{microtype}
\usepackage{xcolor}
\usepackage{graphicx}
\usepackage{multirow}

% \title{BarrierShare: Barrier-Limited Partial Sharing for Multi-Stage Agent Workflows}

\title{Adaptive Agent Selection for Multi-Stage LLM Workflows: 
An Online Learning Approach with Limited Feedback Sharing}

\author{Anonymous Authors}

\newcommand{\towrite}[1]{\noindent\textbf{What to write.} #1}

\begin{document}

\maketitle

\begin{abstract}
% Hierarchical decision systems are a general learning setting in which a terminal outcome is determined by sequential decisions made across multiple stages or agents.
% However, prior work has largely focused on the end-to-end bandit regime, leaving hierarchical learning with partially propagated terminal feedback underexplored. 
% To address this gap, we propose BarrierShare, a distributed online learning framework for hierarchical decision systems with partially propagated terminal feedback.
% Our key insight is that partial propagation is a structural property of the hierarchy: terminal information can be aggregated within shared regions, but must switch to local bandit-style learning once a propagation barrier is reached. 
% Concretely, BarrierShare combines recursive aggregation over barrier-free shared regions, local bandit updates at risky nodes, and barrier-limited delta propagation that stops at the first structural barrier.
% Experiments on simulation and fixed-stage agent workflows evaluate whether these mechanisms improve regret, preserve the predicted long-safe-suffix and depth-scaling behavior, and remain practical under LLM-call, token, and latency costs.

Driven by large language models (LLMs), agentic systems have become increasingly powerful in solving real-world tasks, which can be naturally modeled as multi-stage workflows. 
These tasks often require distinct capabilities at different stages, 
and different LLM agents exhibit heterogeneous abilities. 
To complete a multi-stage task efficiently, 
it is necessary to select the most favorable agent for each stage, making effective agent selection a critical problem.
To address this challenge, we develop an efficient online learning approach that adaptively learns the optimal agent selection for each stage in a distributed manner. 
We first formulate this problem as a multi-stage multi-armed bandit (MAB) problem. Instead of applying existing multi-stage MAB algorithms, 
we observe that overall performance can be significantly improved 
by allowing agents to share only a limited amount of feedback information. 
Leveraging this insight, we propose an EXP3-style algorithm for agent selection. Furthermore, to accommodate scenarios where not all agents are willing to share feedback, we extend the proposed algorithm to a setting where only a subset of agents opt into feedback sharing, and we show that the algorithm can still benefit from partial feedback.
Both synthetic data-driven simulations and real-world LLM agent-based experiments demonstrate that the proposed agent selection mechanism outperforms existing approaches.
\end{abstract}

\section{Introduction (Jinyu)}

Multi-stage agent workflows are becoming a common design pattern for practical language-agent systems.
Many deployed agents do not act as unconstrained planners, but instead operate through structured decision procedures with domain-specific tools, policy constraints, and workflow-defined stages.
This structure appears in prior web-agent and tool-agent settings, although in different forms.
For example, workflow-guided exploration induces high-level workflows from demonstrations to constrain action choices in multi-step web tasks such as booking flights or replying to emails~\cite{liu2018workflow}.
Similarly, $\tau$-bench studies realistic customer-service interactions in which language agents must converse with users, invoke domain-specific API tools, and follow policy guidelines in airline and retail domains~\cite{yao2024tau}.
These works do not prescribe our tree model, but they motivate the view that agent execution is often a structured trajectory rather than a flat action choice.
We abstract such trajectories as fixed-stage workflows, with stages such as user grounding, entity resolution, evidence retrieval, policy interpretation, adjudication, and terminal execution.

In this abstraction, each completed path produces a terminal scalar cost, which can still be used by the selected path through local bandit updates.
The additional question is whether the persistent update induced by this cost can be accumulated into shared ancestor weights.
An externally specified feedback-sharing policy determines whether this induced update is admissible for a shared prefix, or whether it must remain local to a branch.
This distinction matters because treating all induced updates as local wastes reusable statistical signal, while treating all persistent updates as globally shared can expose provider-specific or specialist participation patterns.
This raises a basic question for workflow-constrained agents: when can the update induced by terminal feedback be reused as a shared hierarchical learning signal, and when should it remain local?

A natural formal lens for such systems is multi-stage online learning with terminal-only feedback, where a learner selects a full path through a hierarchy and observes only the realized terminal cost~\cite{hou2024distributed}.
This end-to-end bandit formulation is a robust and conservative starting point: it assumes no more than the cost of the executed path and learns through local bandit estimates along the realized trajectory.
However, it does not model when the same terminal observation may induce a persistent aggregate update that is reusable by other policies inside a feedback-admissible subtree, while being blocked at structural barriers.
\textbf{L1: Structural flattening.}
The realized terminal cost is used only through the sampled trajectory, without distinguishing whether its induced update is also valid for other complete policies sharing a safe prefix or subtree.
\textbf{L2: Missing barrier-limited sharing.}
Existing end-to-end bandit formulations do not capture the intermediate regime where the sampled cost can support local bandit learning, while the persistent shared update is allowed to propagate only inside a barrier-free region.

\emph{When a terminal outcome is admissible for a prefix or subtree under a given feedback-sharing policy, how should a hierarchical learner reuse the induced update without violating structural barriers?}
Our key insight is that feedback reuse should be governed by structural admissibility.
Conceptually, terminal feedback can induce aggregate updates inside policy-admissible safe regions, but propagation must stop at ownership, telemetry, provider, or specialist barriers.
Algorithmically, BarrierShare realizes this principle through an exponential-weight delta.
When a selected leaf is updated, its accumulated score changes, which induces a change in its exponential weight; this induced weight change defines a \emph{delta} that can be added to admissible ancestor aggregate weights.
This shared delta differs from the current-round scalar cost used in a local bandit update: over time, accumulated deltas determine the subtree weights used by ancestor decisions, thereby forming a persistent statistical profile of downstream branches.
When sharing is complete within a subtree, recursive aggregation makes the subtree behave like a single leaf-level adversarial bandit over its leaves, so the regret depends on the number of leaves in that safe subtree rather than accumulating an education penalty across every internal layer.
When a structural barrier is encountered, the shared delta must stop, and the upstream decision must instead rely on local bandit-style estimation.

Such barriers are natural in agent workflows with data-silo, provider, contractual, or telemetry restrictions.
A downstream module may return the current task output and scalar cost, while still prohibiting a globally visible long-run update trace.
This distinction is motivated by analogous leakage risks in collaborative learning, where shared updates can reveal unintended properties of a participant's data or membership information~\cite{melis2019unintended}, and in gradient-based distributed learning, where gradients can even reconstruct private training examples~\cite{zhu2019deep}.
It is also relevant to deployed model services, where repeated black-box queries can expose commercially valuable model behavior or enable model extraction~\cite{tramer2016stealing}.
In sensitive domains such as healthcare, multi-institutional learning is often designed to avoid direct patient-data sharing across silos~\cite{sheller2020federated}.
By analogy, cumulative shared deltas from a private retrieval specialist, clinical-data module, enterprise knowledge component, or third-party provider may reveal the data domain, provider identity, task distribution, or specialization profile behind that branch.
We therefore model these interfaces as externally specified feedback-sharing barriers, not as a formal privacy guarantee.
The relevant regime is neither ``share everything'' nor ``share nothing'': current terminal costs may remain usable through local bandit channels, while persistent shared deltas may propagate upward only until the first structural barrier.
To address this setting, we propose \textbf{BarrierShare}, a hierarchical online learning algorithm that performs shared delta aggregation inside barrier-free regions and switches to local bandit learning at risky or barrier-limited interfaces.
Figure~\ref{fig:overview} illustrates this progression from end-to-end bandit feedback to full-share propagation and finally to the barrier-limited regime studied in this paper.

\textbf{Our contributions are summarized as follows.}
\textbf{(1) A barrier-limited feedback model.}
We formalize an intermediate hierarchical feedback regime in which terminal scalar costs remain usable for local bandit learning, but the persistent shared deltas induced by those costs can propagate only within structurally admissible regions.
This model separates current-round cost feedback from long-run aggregate update traces, and introduces safe regions, structural barriers, and risky depth as the key objects of analysis.
\textbf{(2) The BarrierShare algorithm.}
We propose \textbf{BarrierShare}, a single hierarchical online learning algorithm that combines recursive exponential-weight aggregation inside safe subtrees with local bandit updates at barrier-limited risky interfaces.
A selected shared leaf updates its exponential weight, propagates the induced delta through admissible ancestors, and stops the propagation at the first structural barrier.
\textbf{(3) Risky-depth regret and validation.}
We prove regret bounds governed by the risky skeleton depth \(R\), showing how BarrierShare interpolates between fully shared safe aggregation and depth-dependent bandit learning.
In particular, long safe suffixes reduce the effective depth in the leading bandit term from the full workflow depth \(L\) to the risky depth \(R\), up to safe-subtree aggregation terms.
We validate these predictions through regret-aligned simulations, including long-safe-suffix and depth-scaling studies, and through fixed-stage agent workflow experiments.

\section{Related Work (Jinyu)}

BarrierShare is related to three lines of work, distinguished primarily by their information structure.
First, multi-stage bandits study terminal-only feedback in distributed hierarchies, but use the realized outcome only through local bandit updates.
Second, feedback-graph, semi-bandit, combinatorial bandit, and partial-monitoring models enrich the current-round observation, whereas BarrierShare keeps terminal-only observations and controls the propagation of the induced update.
Third, workflow-constrained agents motivate fixed-stage policies and terminal evaluation, but typically optimize execution within a trajectory rather than repeated online selection among workflow policies.
Table~\ref{tab:feedback_regimes} summarizes these regimes.

\begin{table}[t]
\centering
\small
\caption{Feedback regimes relevant to BarrierShare. Unlike stronger-feedback models, BarrierShare does not enrich the current-round observation; it restricts the scope over which the terminal-cost-induced update trace may be reused.}
\label{tab:feedback_regimes}
\begin{tabular}{lllll}
\toprule
Regime & Observation & Reused object & Reuse scope & Example \\
\midrule
End-to-end bandit & sampled terminal cost & none & none & \citet{hou2024distributed} \\
Complete / one-hop & child/action-level feedback & current-round estimates & local node & Hou complete feedback \\
Full-share endpoint & sampled terminal cost & weight delta & all ancestors & conceptual endpoint \\
BarrierShare & sampled terminal cost & weight delta & safe subtrees & ours \\
\bottomrule
\end{tabular}
\end{table}

\subsection{Multi-Stage Bandits and Distributed Online Learning}

Classical adversarial bandits study a single learner that selects one action and observes only the loss of the selected action, with algorithms such as EXP3 providing no-regret guarantees in the one-stage setting~\citep{auer2002nonstochastic}.
Multi-stage systems introduce a qualitatively different difficulty: the terminal outcome is determined by a sequence of decisions across stages, and upstream decisions affect not only immediate performance but also whether downstream learners receive enough traffic to learn.
This creates the education--exploration--exploitation trilemma formalized by \citet{hou2024distributed}.

The closest theoretical baseline for BarrierShare is Hou's multi-stage bandit model.
In that formulation, a system is represented as a tree: each internal node selects a child, each leaf produces a terminal cost, and the selected path observes only the realized end-to-end outcome.
The $\epsilon$-EXP3 algorithm addresses the education problem by mixing uniform exploration with local exponential-weights updates, yielding regret of order \(O(T^{L/(L+1)})\) for an \(L\)-stage system.
This is the conservative endpoint for our work: terminal feedback is available, but it is used only through local bandit-style updates along the sampled trajectory.

Hou also studies a stronger complete one-hop feedback setting, where each node can observe richer feedback about all of its actions and normalized-EG-type updates become possible.
This setting is fundamentally different from BarrierShare: it enriches current-round observations, whereas our model keeps terminal-only scalar feedback and restricts whether the induced persistent update trace can be shared upward.

The two-stage model of \citet{singla2018learning} is related in that a central forecaster interacts with learning experts and can coordinate their learning by blocking whether the selected expert observes feedback.
However, this feedback blocking controls the learning state of selected experts in a two-level expert-advice problem, whereas BarrierShare studies arbitrary-depth hierarchies in which terminal costs remain locally usable but the induced persistent update delta may propagate only through structurally admissible ancestors.

Our setting is also distinct from adversarial MDP and tree-search formulations.
Although a hierarchy can be viewed as a tree of states, these methods typically assume transition-level state observations, simulator access, value estimates, or centralized coordination.
Unlike adversarial MDP formulations, our learner does not observe intermediate states or rewards along the trajectory, and the regret comparator is a stationary root-to-leaf workflow policy rather than an adaptive Markov policy.
Unlike tree-search methods, our tree is a fixed high-level workflow-policy family rather than a low-level action expansion tree, and learning occurs across repeated terminal-cost observations rather than through within-trajectory search.

\subsection{Feedback Regimes in Hierarchical Learning}

A broad online learning literature studies how regret changes when the learner receives feedback beyond the loss of the selected action.
Full-information learning reveals all action losses; semi-bandit and combinatorial bandit models reveal component-level feedback; partial-monitoring and feedback-graph models specify which side observations become available after each action~\citep{cesa2006prediction, mannor2011bandits, audibert2014regret, alon2015online}.
These settings improve learning by enriching the information available in the current round.

BarrierShare differs in the object being reused.
We do not reveal losses of unselected leaves or child-level costs; the only observed scalar remains the terminal cost of the sampled path.
In a shareable region, this cost induces an exponential-weight delta at the sampled leaf, and this delta updates admissible ancestor aggregates that determine future hierarchical sampling probabilities.
Thus, the question is not which additional losses are observed in the current round, but how far the induced update trace may propagate.

Path and combinatorial bandits also consider actions with internal structure, such as paths, subsets, or other combinatorial objects~\citep{cesa2012combinatorial, audibert2014regret}.
However, semi-bandit variants typically reveal component losses, and bandit variants treat the entire path as a flat arm without providing a mechanism for recursively reusing terminal feedback across safe prefixes.
BarrierShare lies between these extremes: it never observes component-level losses, but it also avoids flattening the hierarchy whenever the induced update is admissible for a safe subtree.

As a useful conceptual endpoint, unrestricted full-share recursive aggregation would allow a selected shareable leaf to update all ancestor aggregates through its weight delta.
BarrierShare retains this mechanism only inside structurally admissible regions; once a barrier is reached, upstream decisions revert to local bandit learning.
Our motivation is also consistent with observations from collaborative and federated learning that shared updates may expose information about data, participants, or model behavior~\citep{melis2019unintended, zhu2019deep, tramer2016stealing, sheller2020federated}.
We do not claim a privacy guarantee: barriers are externally specified sharing constraints, and our goal is to characterize their online-learning consequences.

\subsection{Workflow-Constrained Language Agents}

Recent language-agent systems increasingly structure execution through tools, search, verifiers, memory, or workflow modules.
Tool-using and reasoning-acting agents interleave reasoning, actions, and observations when interacting with external environments~\citep{yao2022react, schick2023toolformer}.
Web-agent and computer-use benchmarks expose long-horizon tasks with tool calls, state tracking, policy constraints, and terminal success criteria~\citep{zhou2023webarena, yao2024tau}.
These works motivate the view that many agent executions are organized around constrained decision points rather than a single flat action.

Importantly, the tree in our formulation is not a low-level action tree over clicks, tokens, or simulator states; it is a fixed high-level policy family whose leaves are complete workflow policies or macro-strategies.
A workflow policy may choose among prompts, tools, specialist modules, API calls, or standard operating procedures at different stages.
Thus, our abstraction is closer to repeated selection among candidate workflow policies than to online search over low-level environment actions.

Recent agent work improves trajectory execution through workflow guidance, process reward models and verifiers, value-guided search, memory retrieval, distillation, or computer-use data engines~\citep{liu2018workflow, chae2025webshepherd, mendes2025selftaught, kang2025agentdistillation, zhang2025gmemory, wang2025opencua}.
These methods change how an agent executes, evaluates, retrieves, or learns within a trajectory.
Unlike prior agent works that improve the executor, verifier, reward model, search procedure, or memory within a trajectory, BarrierShare studies the learning layer above execution: how to repeatedly select among fixed-stage workflow policies when only terminal costs are observed and only part of the induced update trace can be propagated upward.


% \input{files/3-problem}


\section{Problem Formulation (Guocong){\color{red}(V2)}}

% We study hierarchical online learning on a rooted decision tree under three feedback regimes: the conservative end-to-end bandit regime, an idealized full-share endpoint, and the barrier-limited sharing regime that is the focus of this paper. This section defines the game, the feedback semantics, and the regret objective, without introducing the BarrierShare algorithm itself.

In this section, we first formalize the agent selection problem for multi-stage agent workflows. Then, we cast this problem into a multi-stage MAB problem.

\subsection{Agent Selection for Multi-Stage Workflow: A Tree-Structured Decision Process}

\paragraph{Notation.}
Let $\mathcal{T}$ be a rooted tree with root $r$, leaf set $\mathcal{L}$, and child set $C_i$ for each non-leaf node $i$. For any node $i$, let $L(i) \subseteq \mathcal{L}$ denote the set of leaf descendants of $i$. At round $t$, the learner selects a root-to-leaf path $\ell_t \in \mathcal{L}$ by making sequential decisions along the tree. We write $p_t(j \mid i)$ for the probability of selecting child $j \in C_i$ at node $i$, and $\Pi_t(i)$ for the probability of reaching node $i$ under the round-$t$ decision rule. In particular, $\Pi_t(\ell)$ denotes the probability of selecting a leaf $\ell \in \mathcal{L}$. Each leaf $\ell$ incurs a terminal cost $c_t(\ell) \in [0,1]$. For any node $i$, let $Y_\eta[i,T]$ denote the cumulative expected loss of the learner over $T$ rounds starting from $i$, and let $Y_\ast[i,T]$ denote the cumulative loss of the optimal stationary policy starting from $i$.

\textbf{Requirement 1: Privacy-preserved feedback sharing.}

\textbf{Requirement 2: Distributed and independent decision making.}

\subsection{Multi-Stage MAB Formulation}

We model a multi-stage decision system as a rooted tree $\mathcal{T}$ of depth $L$. Each internal node represents a decision point, and each leaf represents a complete root-to-leaf decision path. At each round $t = 1,\ldots,T$, the learner starts at the root and sequentially selects one child at each internal node until reaching a leaf $\ell_t \in \mathcal{L}$. Equivalently, the round-$t$ decision can be written as a path
\[
\ell_t = (a_{1,t}, \ldots, a_{L,t}),
\]
where $a_{k,t}$ is the stage-$k$ choice along the selected path.

The decision rule is specified by local conditional probabilities $p_t(j \mid i)$ at internal nodes. The induced probability of reaching a node is defined recursively by
\[
\Pi_t(r) = 1,
\qquad
\Pi_t(j) = \Pi_t(i)\, p_t(j \mid i)
\quad \text{for } j \in C_i.
\]
Hence, for a leaf $\ell$, $\Pi_t(\ell)$ is the probability of selecting the entire root-to-leaf path ending at $\ell$.

After selecting $\ell_t$, the learner incurs and observes the terminal cost $c_t(\ell_t)$. We compare the learner against the best stationary policy that fixes, at every internal node, a single child choice throughout all rounds. Such a policy induces a fixed leaf in every subtree. Define recursively
\[
Y_\ast[i,T] =
\begin{cases}
\sum_{t=1}^T c_t(i), & i \in \mathcal{L},\\[4pt]
\min_{j \in C_i} Y_\ast[j,T], & i \notin \mathcal{L}.
\end{cases}
\]
The cumulative expected loss of the learner from node $i$ is
\[
Y_\eta[i,T] := \sum_{t=1}^T \mathbb{E}[y_\eta[i,t]],
\]
where $y_\eta[i,t]$ is the random loss incurred at round $t$ if the learner starts from $i$ and follows its round-$t$ decision rule thereafter. Our root regret objective is
\[
\mathrm{Reg}_T := Y_\eta[r,T] - Y_\ast[r,T].
\]

This formulation provides the formal backbone of multi-stage end-to-end learning before any additional feedback structure is introduced.

\subsection{End-to-End Bandit Regime}

We first consider the conservative regime in which the learner only observes the scalar terminal cost of the selected path. Formally, after choosing $\ell_t$, the learner observes only
\[
c_t(\ell_t),
\]
and receives no direct information about the terminal costs of unselected leaves or the intermediate stage-wise contributions along the hierarchy.

Under this regime, terminal feedback is usable only through bandit-style estimators based on the sampling probability of the selected path or selected local decisions. This model is safe in the sense that it requires no assumption beyond observability of the realized end-to-end outcome. At the same time, it is conservative: even when part of the terminal signal could serve as reusable supervision for an upstream shared structure, the model treats the entire signal as path-local. Representative baselines in this regime include EXP3-type and $\epsilon$-EXP3-type methods, but this subsection defines only the feedback model rather than any particular algorithm.

\subsection{Full-Share Delta Regime}

We next define an idealized full-share endpoint in which terminal feedback from a selected leaf can be converted into a sharable delta and propagated upward through the hierarchy.

Let $\mathcal{L}^{\mathrm{shr}} \subseteq \mathcal{L}$ denote the set of leaves whose terminal feedback is fully shareable with their ancestors. When a selected leaf $\ell_t \in \mathcal{L}^{\mathrm{shr}}$ is observed, its terminal feedback induces a sharable delta, denoted by $\Delta_t(\ell_t)$, which can be transmitted upward to all ancestors that fully observe the corresponding shared subtree.

The only invariant we require from this regime is the aggregate-visibility property: for any node $i$ that fully observes a shared subtree, its aggregate shared weight equals the sum of the visible descendant shared leaf weights. That is, if $W_t[i]$ denotes the aggregate weight at node $i$, then
\[
W_t[i] = \sum_{\ell \in L(i) \cap \mathcal{L}^{\mathrm{shr}}} w_t(\ell),
\]
where $w_t(\ell)$ is the shared leaf weight associated with $\ell$. This full-share regime serves as an idealized endpoint that reveals the value of reusable hierarchical feedback, but it does not yet capture structural restrictions on how far such feedback may propagate.

\subsection{Barrier-Limited Sharing Regime}

We now define the barrier-limited regime studied in this paper. Each non-leaf node $i$ is equipped with a gate
\[
g(i) \in \{\mathrm{share}, \mathrm{unshare}\},
\]
which determines whether shared feedback may propagate upward through that node. As before, let $\mathcal{L}^{\mathrm{shr}} \subseteq \mathcal{L}$ denote the set of leaves whose terminal feedback is eligible to initiate shared upload.

At round $t$, if the selected leaf $\ell_t$ belongs to $\mathcal{L}^{\mathrm{shr}}$, then its terminal feedback may initiate an upward shared upload. However, this propagation continues only until the first ancestor whose gate is $\mathrm{unshare}$. Such a node acts as a structural barrier, beyond which the shared delta cannot propagate.

A non-leaf node $i$ is called \emph{safe} if the entire relevant subtree beneath $i$ is fully shareable and fully visible at $i$, namely if
\[
L(i) \subseteq \mathcal{L}^{\mathrm{shr}}
\quad \text{and} \quad
g(u)=\mathrm{share}, \ \forall \text{ non-leaf } u \in \mathrm{subtree}(i).
\]
Otherwise, $i$ is called \emph{risky}. Intuitively, a safe node sits above a barrier-free fully shared region, whereas a risky node lies on or above a region where upward propagation may be blocked.

To quantify the amount of non-shareable structure, define the risky depth recursively by
\[
R(i)=
\begin{cases}
0, & \text{if } i \text{ is safe},\\
1, & \text{if } i \text{ is risky and } C_i \subseteq \mathcal{L},\\
1+\max_{j \in C_i} R(j), & \text{if } i \text{ is risky and } C_i \not\subseteq \mathcal{L}.
\end{cases}
\]
The risky depth of the whole tree is then
\[
R := R(r).
\]
This quantity captures the maximum number of risky decision points encountered on any root-to-leaf path before entering a barrier-free fully shared suffix.

\paragraph{Core question.}
\emph{How can a learner achieve low regret when terminal feedback can be propagated upward only until structural barriers are encountered?}
\emph{How can a learner achieve low regret when terminal feedback can be propagated upward only until structural barriers are encountered?}

\section{Insightful Observation: Large Gain from Limited Information Sharing (Jinyu)}

\section{BarrierShare (Guocong)}
\label{sec:barriershare}

BarrierShare learns how to route probability mass through the tree, how to reuse terminal feedback when sharing is permitted, and how to block propagation when a structural barrier is encountered. The algorithm interpolates between two endpoints: recursive sharing in fully shareable regions and local bandit learning in risky regions.

\subsection{Method Overview}
\label{sec:barriershare_overview}

BarrierShare maintains two coupled mechanisms. The first is \emph{recursive selection}: at each internal node, the learner selects a child either according to aggregated subtree mass or according to a local exploration-exploitation rule. The second is \emph{barrier-limited update}: after observing terminal feedback at the selected leaf, the learner performs local updates on the sampled risky path; if the selected leaf is shareable, it additionally propagates the induced leaf delta upward through the shareable prefix of the path.

Safe regions reuse feedback through exact aggregation over shared leaves, whereas risky regions only admit local credit assignment from the sampled trajectory. These two mechanisms operate on different state variables but are coupled through the sampled path.

\subsection{Recursive Selection}
\label{sec:barriershare_selection}

Starting from the root, BarrierShare recursively selects one child at each visited internal node until a leaf is reached.

\paragraph{Safe nodes.}
For any safe node $u$, child selection is induced by aggregated subtree weights:
\begin{equation}
    p_t(v\mid u) = \frac{W_t(v)}{W_t(u)},
    \qquad v \in \mathsf{ch}(u).
    \label{eq:safe_selection}
\end{equation}
Here $W_t(u)$ is the exact aggregate of shared leaf weights in the shareable subtree rooted at $u$. Because safe subtrees are closed under descendants, the same rule applies recursively to every child subtree below a safe node.

\paragraph{Risky nodes.}
If the current node is risky, BarrierShare uses an $\epsilon$-greedy exploration rule:
\begin{equation}
    p_t(v\mid u)
    =
    \begin{cases}
    \dfrac{1}{|\mathsf{ch}(u)|}, & \text{with probability } \epsilon_i, \\[8pt]
    \dfrac{\exp(\eta \theta_t(v\mid u))}{\sum_{v' \in \mathsf{ch}(u)} \exp(\eta \theta_t(v'\mid u))}, & \text{with probability } 1-\epsilon_i,
    \end{cases}
    \qquad v \in \mathsf{ch}(u),
    \label{eq:risky_selection}
\end{equation}
where $\epsilon_i \in (0,1)$ controls exploration at node $u$. This rule means that with probability $\epsilon_i$ the learner samples uniformly at random, and with probability $1-\epsilon_i$ it samples according to the exponential-weight distribution.

The probability of the full sampled path is
\begin{equation}
    P_t(\pi_t) \coloneqq \prod_{k=0}^{D_t-1} p_t(u_{k+1}\mid u_k).
    \label{eq:path_prob}
\end{equation}
We use $P_t(\pi_t)$ only to describe the realized trajectory probability; the node-local update below uses the reach probability $\rho_t(u)$ together with the local choice probability $p_t(v\mid u)$.

\begin{algorithm}[t]
\caption{\textsc{BarrierShare}}
\label{alg:barriershare}
\KwIn{Tree $\mathcal{T}$, shared leaf set $\mathcal{L}_{\mathrm{shr}}$, barrier mask $g$, learning rate $\eta$, exploration rate $\epsilon$}
\KwOut{Sampled leaf $\ell_t$ and updated states}

Initialize $\phi_1(\ell)=0$ and $w_1(\ell)=\exp(-\eta \phi_1(\ell))$ for all $\ell\in\mathcal{L}_{\mathrm{shr}}$\;
Initialize $W_1(u)=\sum_{\ell\in\mathcal{L}(u)\cap \mathcal{L}_{\mathrm{shr}}} w_1(\ell)$ for all safe nodes $u$\;
Initialize $\theta_1(v\mid u)=0$ for all directed edges $(u,v)$\;

Set $u \leftarrow r$\;
\While{$u \notin \mathcal{L}$}{
    \eIf{$u$ is safe}{
        sample $v \in \mathsf{ch}(u)$ according to \eqref{eq:safe_selection}\;
    }{
        sample $v \in \mathsf{ch}(u)$ according to \eqref{eq:risky_selection}\;
    }
    record selected edge $(u,v)$ and set $u \leftarrow v$\;
}
Let the sampled leaf be $\ell_t \leftarrow u$ and observe terminal cost $c_t$\;

\ForEach{risky node $u$ on $\pi_t$ with sampled child $v$}{
    compute
    \[
        \widehat{c}_t(u,v)
        \coloneqq
        \frac{c_t \cdot \mathbf{1}\{(u,v)\in\pi_t\}}{\rho_t(u)\,p_t(v\mid u)}\;,
    \]
    and update
    \[
        \theta_{t+1}(v\mid u) = \theta_t(v\mid u) - \eta \widehat{c}_t(u,v)\;.
    \]
}

\eIf{$\ell_t \in \mathcal{L}_{\mathrm{shr}}$}{
    update $\phi_{t+1}(\ell_t)=\phi_t(\ell_t)+c_t$ and recompute $w_{t+1}(\ell_t)=\exp(-\eta\phi_{t+1}(\ell_t))$\;
    compute the induced leaf delta $\Delta_t(\ell_t)=w_{t+1}(\ell_t)-w_t(\ell_t)$\;
    additionally propagate $\Delta_t(\ell_t)$ upward along ancestors until the first node with $g(u)=\mathrm{unshare}$, and stop there\;
}{
    do not generate any reusable ancestor update\;
}
\end{algorithm}

\subsection{Risky Local Update}
\label{sec:barriershare_risky_update}

The risky update is a node-local bandit update. Let $u$ be a risky node on the sampled path and let $v$ be the chosen child of $u$. Since the learner reaches $u$ with probability $\rho_t(u)$ and then chooses $v$ with probability $p_t(v\mid u)$, we use the importance-weighted estimator
\begin{equation}
    \widehat{c}_t(u,v)
    \coloneqq
    \frac{c_t \cdot \mathbf{1}\{(u,v)\in\pi_t\}}{\rho_t(u)\,p_t(v\mid u)}.
    \label{eq:iw_estimator}
\end{equation}
The local score update is
\begin{equation}
    \theta_{t+1}(v\mid u)
    =
    \theta_t(v\mid u) - \eta \widehat{c}_t(u,v).
    \label{eq:risky_local_update}
\end{equation}

This update only affects the visited risky edge. It does not modify subtree weights above the barrier, and it does not assume that the terminal cost is reusable outside the sampled local interface.

\subsection{Shared Leaf Update and Delta Propagation}
\label{sec:barriershare_shared_update}

If the selected leaf $\ell_t$ is shareable, then the terminal feedback is also eligible for reusable aggregation. We update the leaf score by
\begin{equation}
    \phi_{t+1}(\ell_t) = \phi_t(\ell_t) + c_t,
    \label{eq:leaf_score_update}
\end{equation}
which induces the updated leaf weight
\begin{equation}
    w_{t+1}(\ell_t) = \exp\!\bigl(-\eta \phi_{t+1}(\ell_t)\bigr).
    \label{eq:leaf_weight_update}
\end{equation}
We then define the induced leaf delta as
\begin{equation}
    \Delta_t(\ell_t) \coloneqq w_{t+1}(\ell_t) - w_t(\ell_t).
    \label{eq:delta_def}
\end{equation}

For every ancestor $u$ of $\ell_t$ such that the path from $u$ to $\ell_t$ contains only share-labeled nodes, we propagate $\Delta_t(\ell_t)$ upward and refresh the corresponding subtree weights. The propagation stops at the first barrier node, and no ancestor above that node receives the reusable update.

This rule keeps multiplicative leaf updates and additive subtree aggregation consistent: leaf scores are updated through the exponential map, while subtree weights are maintained as exact sums of shared leaf weights. In particular, for every ancestor that receives the propagated delta, the updated subtree weight remains equal to the sum of the updated shared descendant leaf weights.

If $\ell_t \notin \mathcal{L}_{\mathrm{shr}}$, then no reusable ancestor update is generated. The terminal cost still contributes to the risky local update on the sampled path, but it is not propagated into any subtree weight.

\subsection{Endpoint Reductions}
\label{sec:barriershare_reductions}

\paragraph{All-share regime.}
If every node is shareable and every leaf is shared, then every internal node is safe. In this case, BarrierShare reduces to a fully recursive sharing mechanism in which all routing probabilities are induced by subtree weights and every terminal feedback update propagates to the root.

\paragraph{All-unshare regime.}
If no reusable sharing is available, then every non-root node is risky. In this case, BarrierShare reduces to a purely local hierarchical bandit learner that updates only the sampled risky edges along the path.

\paragraph{Intermediate regime.}
In the general case, safe regions reuse shared feedback through exact subtree aggregation, while risky regions remain local and are updated only through bandit feedback. This is the regime of primary interest, because it interpolates between full sharing and no sharing in a structurally interpretable way.

\subsection{Application to Hierarchical Agent Workflows}
\label{sec:barriershare_agent_instantiation}

The same formalism naturally instantiates hierarchical agent workflows. In that setting, each internal node corresponds to a stage-level routing decision and each leaf corresponds to a complete terminal macro-policy. The learner selects which leaf policy to execute, while the executor carries out the leaf-level trajectory and returns a single terminal cost after completion.

This mapping separates routing from execution. BarrierShare learns how to route probability mass through the tree, but it does not require step-wise supervision inside each leaf. Therefore, the same algorithm can be used for synthetic tree simulations and for real hierarchical agent pipelines, provided that the terminal cost is observed only after the selected leaf terminates.

\section{Experiments}
\label{sec:experiments}

Our experiments evaluate whether barrier-limited feedback sharing improves learning in fixed-stage agent workflows when only terminal execution feedback is observed. The central question is not whether a single workflow can be hand-tuned to perform well, but whether a learner that respects structural sharing barriers can identify cheaper post-switch workflow policies more reliably than endpoint sharing regimes, flat or local bandit baselines, and LLM-only routing baselines.

We focus on a telecom MMS diagnosis and repair workflow. Each episode selects a complete five-stage root-to-leaf workflow policy, executes it with the same LLM executor, and observes a scalar terminal cost. The tree learner receives only this terminal outcome for update. It does not receive stage-wise supervision, oracle labels during learning, or direct feedback for unchosen leaves.

The final experimental suite uses three profile-switch tree families:
\textbf{v4}, a compact binary mixed tree for mechanism isolation;
\textbf{v6}, a 30-leaf tree that keeps the trap-switch structure while increasing post-switch exploration pressure; and
\textbf{v3}, a 250-leaf efficient-anchor tree that tests the same mechanism at a larger workflow-policy scale. Across all three trees, we compare the same set of eleven methods.

\subsection{Questions}
\label{sec:exp_questions}

The experiments are organized around four questions.

\textbf{(Q1) Does partial sharing help?}
We compare the main 4/5-share BarrierShare configuration against all-share and all-unshare endpoints, a lower-sharing 2/5 variant, and standard bandit baselines.

\textbf{(Q2) Is the effect robust across tree scale?}
We test the same method set on the compact v4 tree, the intermediate v6 tree, and the larger v3 250-leaf tree.

\textbf{(Q3) Does structured learning beat unstructured or LLM-only routing?}
We compare BarrierShare not only with EXP-style algorithms, but also with naive mixing, random routing, theta-guided LLM routing, and agent-only LLM routing.

\textbf{(Q4) Are improvements due to valid shared-basin reuse rather than unsafe transfer?}
We report terminal-quality and execution diagnostics separately from aggregate cost so that conservative transfer behavior cannot be mistaken for successful local repair.

\subsection{Workflow Environment}
\label{sec:exp_workflow_env}

The environment is a fixed five-stage telecom MMS workflow. A root-to-leaf path specifies one complete workflow policy:
(1) user grounding,
(2) customer-line resolution,
(3) observed feature extraction,
(4) blocker adjudication and repair planning, and
(5) terminal execution decision.

The executor is shared across methods. In the current experimental configuration, the LLM execution layer uses the Stage~4/5 contract prompt family, prompt version \texttt{v11b}, the clean-v3 execution setup, and \texttt{gpt-4.1-mini} as the workflow model. The executor returns a terminal outcome and resource diagnostics for each selected leaf. The learner then updates from the resulting scalar cost.

We use a profile-switch schedule. Early episodes contain trap or pre-switch pressure, while later episodes favor deep target/shared-basin workflows. This schedule is intended to expose the main failure mode of unsafe or overly conservative methods: a method can look good early by exploiting fast/trap behavior, but it must eventually discover cheaper post-switch shared-basin policies.

\subsection{Tree Families}
\label{sec:exp_tree_families}

Table~\ref{tab:tree_families} summarizes the three tree families used in the final experiments. All trees have five workflow stages. They differ in the number of leaves, the amount of shared structure, and the difficulty of post-switch exploration.

\begin{table}[t]
    \centering
    \small
    \caption{Tree families used in the final experimental suite. All families use the same five-stage telecom MMS executor, but differ in topology and sharing pattern.}
    \label{tab:tree_families}
    \begin{tabular}{llp{7.1cm}}
        \toprule
        Family & Size & Purpose \\
        \midrule
        v4 binary mixed & 8 leaves & Compact mechanism test with fast trap leaves and a shared deep target basin. Used to isolate whether the sharing rule behaves differently from EXP-style and endpoint regimes. \\
        v6 small30 & 30 leaves & Intermediate profile-switch tree with 4/5 sharing, six trap leaves, and twenty-four deep target leaves. Used to test whether the mechanism remains useful when post-switch exploration is less trivial. \\
        v3 efficient-anchor & 250 leaves & Larger efficient-anchor profile-switch tree with 4/5 sharing. Used as the main scale test for the final comparison. \\
        \bottomrule
    \end{tabular}
\end{table}

For each family we instantiate the same share variants whenever the topology supports them:
all-share, 2/5-share, 4/5-share, and all-unshare. The topology, executor, prompt contract, and terminal-cost definition are held fixed within a family; only the gate map used by the BarrierShare-style learner changes.

\subsection{Methods}
\label{sec:exp_methods}

Each tree is evaluated with eleven methods. We group them by what information they use.

\paragraph{BarrierShare share variants.}
The main method is \texttt{ps\_4/5share}, which uses the intended 4/5 gate map for the tree. We also run \texttt{ps\_allshare}, \texttt{ps\_2/5share}, and \texttt{ps\_allunshare}. These are not merely extra baselines; they are structural controls. \texttt{ps\_allshare} tests the optimistic endpoint in which feedback can propagate everywhere, \texttt{ps\_allunshare} tests the conservative endpoint in which reusable propagation is disabled, and \texttt{ps\_2/5share} tests whether the method degrades gracefully under reduced sharing.

\paragraph{Bandit baselines.}
\texttt{exp} denotes the direct multistage EXP baseline. \texttt{eps} denotes the epsilon-EXP baseline with explicit random exploration. \texttt{exp\_local} denotes the local multistage EXP baseline that keeps updates more local than the direct shared variant.

\paragraph{Heuristic and unstructured baselines.}
\texttt{naivemix} mixes candidate workflow policies without the BarrierShare propagation rule. \texttt{random} samples legal paths without learning and serves as a sanity check.

\paragraph{LLM routing baselines.}
\texttt{theta guided} exposes algorithmic stage-local scores to an LLM selector, which then chooses the next stage action. \texttt{agent\_only} asks the LLM to choose stage actions without algorithmic theta guidance. These baselines test whether the learned tree policy adds value beyond prompting an agent to route through the workflow.

\begin{table}[t]
    \centering
    \small
    \caption{Methods compared on each tree family.}
    \label{tab:methods}
    \begin{tabular}{lll}
        \toprule
        Paper name & Implementation name & Role \\
        \midrule
        \texttt{ps\_allshare} & PS on all-share gate map & Optimistic sharing endpoint \\
        \texttt{ps\_2/5share} & PS on 2/5 gate map & Reduced-sharing control \\
        \texttt{ps\_4/5share} & PS on 4/5 gate map & Main BarrierShare setting \\
        \texttt{ps\_allunshare} & PS on all-unshare gate map & No-sharing endpoint \\
        \texttt{eps} & \texttt{epsilon\_exp3} & Exploration baseline \\
        \texttt{exp} & \texttt{direct\_multistage\_exp3} & Direct multistage EXP baseline \\
        \texttt{exp\_local} & \texttt{direct\_multistage\_exp3\_local} & Local multistage EXP baseline \\
        \texttt{naivemix} & \texttt{naive\_mixed} or \texttt{naive\_mixed\_avg} & Heuristic mixture baseline \\
        \texttt{theta guided} & \texttt{theta\_guided\_agent} & LLM selector with theta guidance \\
        \texttt{random} & \texttt{random\_path} & Random-path sanity check \\
        \texttt{agent\_only} & \texttt{agent\_only} & LLM-only routing baseline \\
        \bottomrule
    \end{tabular}
\end{table}

\subsection{Evaluation Protocol}
\label{sec:exp_protocol}

Within each tree family, all methods use the same task schedule, executor, model, prompt family, cost function, and random-seed convention. The current final runs use three 100-episode repetitions per method-tree condition. Unless otherwise stated, learning hyperparameters are fixed at $\eta=0.3$ and $\epsilon=0.01$ for the relevant EXP-style and BarrierShare-style learners, with the shared-update scale fixed by the corresponding PS share-variant configuration.

The primary metric is average total cost over episodes. We also report post-switch cost, terminal penalty, exact-match or clean-success rate when available, and resource diagnostics such as LLM call count, token usage, API cost, and wall-clock time. Because profile-switch experiments can be distorted by early trap episodes, we report aggregate metrics and post-switch metrics separately.

For mechanism checks, we log the selected path, stage choices, gate pattern, update type, shared-update activity, and whether execution reached local repair, subset repair, or transfer. These diagnostics are used to interpret the results but are not used to change the policy during evaluation.

\subsection{Main Results}
\label{sec:exp_main_results}

Table~\ref{tab:main_results} is the main result table. It will report the same eleven methods on v4, v6, and v3. Lower cost is better; higher success is better.

\begin{table*}[t]
    \centering
    \small
    \caption{Main workflow results. Results and analysis are intentionally left blank until the final runs are complete.}
    \label{tab:main_results}
    \begin{tabular}{llcccccc}
        \toprule
        Tree & Method & Total cost $\downarrow$ & Post-switch cost $\downarrow$ & Terminal penalty $\downarrow$ & Success $\uparrow$ & Calls $\downarrow$ & Tokens $\downarrow$ \\
        \midrule
        v4 & \texttt{ps\_allshare} & -- & -- & -- & -- & -- & -- \\
        v4 & \texttt{ps\_2/5share} & -- & -- & -- & -- & -- & -- \\
        v4 & \texttt{ps\_4/5share} & -- & -- & -- & -- & -- & -- \\
        v4 & \texttt{ps\_allunshare} & -- & -- & -- & -- & -- & -- \\
        v4 & \texttt{eps} & -- & -- & -- & -- & -- & -- \\
        v4 & \texttt{exp} & -- & -- & -- & -- & -- & -- \\
        v4 & \texttt{exp\_local} & -- & -- & -- & -- & -- & -- \\
        v4 & \texttt{naivemix} & -- & -- & -- & -- & -- & -- \\
        v4 & \texttt{theta guided} & -- & -- & -- & -- & -- & -- \\
        v4 & \texttt{random} & -- & -- & -- & -- & -- & -- \\
        v4 & \texttt{agent\_only} & -- & -- & -- & -- & -- & -- \\
        \midrule
        v6 & \texttt{ps\_allshare} & -- & -- & -- & -- & -- & -- \\
        v6 & \texttt{ps\_2/5share} & -- & -- & -- & -- & -- & -- \\
        v6 & \texttt{ps\_4/5share} & -- & -- & -- & -- & -- & -- \\
        v6 & \texttt{ps\_allunshare} & -- & -- & -- & -- & -- & -- \\
        v6 & \texttt{eps} & -- & -- & -- & -- & -- & -- \\
        v6 & \texttt{exp} & -- & -- & -- & -- & -- & -- \\
        v6 & \texttt{exp\_local} & -- & -- & -- & -- & -- & -- \\
        v6 & \texttt{naivemix} & -- & -- & -- & -- & -- & -- \\
        v6 & \texttt{theta guided} & -- & -- & -- & -- & -- & -- \\
        v6 & \texttt{random} & -- & -- & -- & -- & -- & -- \\
        v6 & \texttt{agent\_only} & -- & -- & -- & -- & -- & -- \\
        \midrule
        v3 & \texttt{ps\_allshare} & -- & -- & -- & -- & -- & -- \\
        v3 & \texttt{ps\_2/5share} & -- & -- & -- & -- & -- & -- \\
        v3 & \texttt{ps\_4/5share} & -- & -- & -- & -- & -- & -- \\
        v3 & \texttt{ps\_allunshare} & -- & -- & -- & -- & -- & -- \\
        v3 & \texttt{eps} & -- & -- & -- & -- & -- & -- \\
        v3 & \texttt{exp} & -- & -- & -- & -- & -- & -- \\
        v3 & \texttt{exp\_local} & -- & -- & -- & -- & -- & -- \\
        v3 & \texttt{naivemix} & -- & -- & -- & -- & -- & -- \\
        v3 & \texttt{theta guided} & -- & -- & -- & -- & -- & -- \\
        v3 & \texttt{random} & -- & -- & -- & -- & -- & -- \\
        v3 & \texttt{agent\_only} & -- & -- & -- & -- & -- & -- \\
        \bottomrule
    \end{tabular}
\end{table*}

\textcolor{gray}{[TODO: Fill Table~\ref{tab:main_results} after the final v4, v6, and v3 runs are complete.]}

\textcolor{gray}{[TODO: Add the main result interpretation here. The analysis should first compare \texttt{ps\_4/5share} against \texttt{exp}, \texttt{eps}, and \texttt{exp\_local}; then compare the PS share variants; then discuss the LLM routing baselines.]}

\subsection{Share-Variant Analysis}
\label{sec:exp_share_variant_analysis}

The share-variant comparison is the most direct test of the paper's mechanism. Holding the tree topology and executor fixed, we vary only the gate map used by the PS learner.

\begin{table}[t]
    \centering
    \small
    \caption{Share-variant comparison.}
    \label{tab:share_variant_results}
    \begin{tabular}{llcccc}
        \toprule
        Tree & PS variant & Total cost $\downarrow$ & Post-switch cost $\downarrow$ & Terminal penalty $\downarrow$ & Shared-update diagnostics \\
        \midrule
        v4 & all-share & -- & -- & -- & -- \\
        v4 & 2/5-share & -- & -- & -- & -- \\
        v4 & 4/5-share & -- & -- & -- & -- \\
        v4 & all-unshare & -- & -- & -- & -- \\
        v6 & all-share & -- & -- & -- & -- \\
        v6 & 2/5-share & -- & -- & -- & -- \\
        v6 & 4/5-share & -- & -- & -- & -- \\
        v6 & all-unshare & -- & -- & -- & -- \\
        v3 & all-share & -- & -- & -- & -- \\
        v3 & 2/5-share & -- & -- & -- & -- \\
        v3 & 4/5-share & -- & -- & -- & -- \\
        v3 & all-unshare & -- & -- & -- & -- \\
        \bottomrule
    \end{tabular}
\end{table}

\textcolor{gray}{[TODO: Add analysis of whether the 4/5-share configuration improves over all-unshare without inheriting the failure modes of all-share. Also report whether 2/5-share behaves as a lower-sharing intermediate point.]}

\subsection{Execution-Quality Analysis}
\label{sec:exp_execution_quality}

Aggregate cost alone is not sufficient in this domain. A method can reduce cost by transferring conservatively or by over-selecting fast paths that avoid difficult local repair. We therefore analyze terminal behavior by execution category:
\texttt{repair\_all},
\texttt{repair\_subset},
and \texttt{transfer}. When available, we also separate local non-transfer tasks from hard hybrid or nonlocal transfer-control tasks.

\begin{table}[t]
    \centering
    \small
    \caption{Execution-quality breakdown.}
    \label{tab:execution_quality}
    \begin{tabular}{llcccc}
        \toprule
        Tree & Method & Local success $\uparrow$ & Transfer rate & Subset quality $\uparrow$ & Hard-transfer safety \\
        \midrule
        v4 & \texttt{ps\_4/5share} & -- & -- & -- & -- \\
        v4 & Best non-PS baseline & -- & -- & -- & -- \\
        v6 & \texttt{ps\_4/5share} & -- & -- & -- & -- \\
        v6 & Best non-PS baseline & -- & -- & -- & -- \\
        v3 & \texttt{ps\_4/5share} & -- & -- & -- & -- \\
        v3 & Best non-PS baseline & -- & -- & -- & -- \\
        \bottomrule
    \end{tabular}
\end{table}

\textcolor{gray}{[TODO: Fill the execution-quality table. The key point to check is whether low total cost corresponds to better local repair quality rather than conservative transfer.]}

\subsection{Learning Dynamics}
\label{sec:exp_learning_dynamics}

For each tree, we plot cumulative cost and moving-average episode cost over the 100-episode horizon. We also mark the switch point so that pre-switch exploration and post-switch target adaptation can be read separately.

\textcolor{gray}{[TODO: Add learning-curve figures for v4, v6, and v3. Use the same method colors across all plots. Include at least \texttt{ps\_4/5share}, \texttt{ps\_allshare}, \texttt{ps\_allunshare}, \texttt{exp}, \texttt{eps}, \texttt{exp\_local}, and the two LLM routing baselines.]}

\subsection{Resource Use}
\label{sec:exp_resource_use}

The executor logs LLM call count, prompt tokens, completion tokens, total tokens, estimated API cost, generation time, and wall-clock episode time. We report these as secondary diagnostics. The resource table is not used to override the terminal-cost metric, but it helps distinguish learning improvements from prompt-length or retry artifacts.

\begin{table}[t]
    \centering
    \small
    \caption{Resource diagnostics.}
    \label{tab:resource_results}
    \begin{tabular}{llcccc}
        \toprule
        Tree & Method & Calls $\downarrow$ & Tokens $\downarrow$ & API cost $\downarrow$ & Wall time $\downarrow$ \\
        \midrule
        v4 & \texttt{ps\_4/5share} & -- & -- & -- & -- \\
        v4 & Best baseline & -- & -- & -- & -- \\
        v6 & \texttt{ps\_4/5share} & -- & -- & -- & -- \\
        v6 & Best baseline & -- & -- & -- & -- \\
        v3 & \texttt{ps\_4/5share} & -- & -- & -- & -- \\
        v3 & Best baseline & -- & -- & -- & -- \\
        \bottomrule
    \end{tabular}
\end{table}

\textcolor{gray}{[TODO: Fill resource diagnostics after final aggregation. If token or API cost differences are small, state that the main effect is terminal quality/cost rather than resource use. If they are large, discuss them separately.]}

\subsection{Planned Interpretation}
\label{sec:exp_planned_interpretation}

The final analysis will distinguish three possible outcomes.

First, if \texttt{ps\_4/5share} consistently improves over both all-share and all-unshare, the result supports the paper's main claim that partial structural sharing is useful beyond endpoint regimes.

Second, if PS variants help on v4 but not on v6 or v3, the result suggests that the mechanism works in small controlled trees but remains sensitive to exploration scale or execution noise.

Third, if EXP-style or LLM-only routing baselines match or outperform the PS variants across trees, the result weakens the empirical claim and the analysis should focus on whether the failure comes from update instability, insufficient shared-basin signal, or executor terminal noise.

\textcolor{gray}{[TODO: Replace this planned interpretation with the actual result-dependent interpretation once the final tables and learning curves are available.]}


\section{Conclusion (Guocong)}

This paper studies hierarchical online learning for multi-stage agent workflows under structured feedback sharing. We begin from the conservative end-to-end bandit view, show that full-share delta propagation reveals the value of reusable upward feedback, and then identify structural barriers as the key reason why real systems lie between full sharing and purely path-local learning. In response, we propose BarrierShare, a unified intermediate algorithm that combines recursive shared aggregation in barrier-free regions with local bandit learning at risky decision points. Our theoretical results show that the resulting learning difficulty is governed by the risky skeleton depth rather than only by the total hierarchy depth, and our simulation and fixed-stage agent experiments provide consistent evidence for this structural picture. A current limitation is that our formulation is tailored to fixed-stage or fixed-for-the-experiment hierarchical workflows, where barriers are modeled as structural feedback-sharing constraints rather than formal privacy guarantees.

% \begin{abstract}
% Hierarchical decision systems are a general learning setting in which a terminal outcome is determined by sequential decisions made across multiple stages or agents.
% However, prior work has largely focused on the end-to-end bandit regime, leaving hierarchical learning with partially propagated terminal feedback underexplored. 
% To address this gap, we propose BarrierShare, a distributed online learning framework for hierarchical decision systems with partially propagated terminal feedback.
% Our key insight is that partial propagation is a structural property of the hierarchy: terminal information can be aggregated within shared regions, but must switch to local bandit-style learning once a propagation barrier is reached. 
% Concretely, BarrierShare combines recursive aggregation over barrier-free shared regions, local bandit updates at risky nodes, and barrier-limited delta propagation that stops at the first structural barrier.
% Experiments on simulation and fixed-stage agent workflows evaluate whether these mechanisms improve regret, preserve the predicted long-safe-suffix and depth-scaling behavior, and remain practical under LLM-call, token, and latency costs.
% \end{abstract}

% \section{Introduction}

% Multi-stage agent workflows are becoming a common design pattern for practical language-agent systems.
% Many deployed agents do not act as unconstrained planners, but instead operate through structured decision procedures with domain-specific tools, policy constraints, and workflow-defined stages.
% This structure appears in prior web-agent and tool-agent settings, although in different forms.
% For example, workflow-guided exploration induces high-level workflows from demonstrations to constrain action choices in multi-step web tasks such as booking flights or replying to emails~\cite{liu2018workflow}.
% Similarly, $\tau$-bench studies realistic customer-service interactions in which language agents must converse with users, invoke domain-specific API tools, and follow policy guidelines in airline and retail domains~\cite{yao2024tau}.
% These works do not prescribe our tree model, but they motivate the view that agent execution is often a structured trajectory rather than a flat action choice.
% We abstract such trajectories as fixed-stage workflows, with stages such as user grounding, entity resolution, evidence retrieval, policy interpretation, adjudication, and terminal execution.

% In this abstraction, each completed path produces a terminal scalar cost, which can still be used by the selected path through local bandit updates.
% The additional question is whether the persistent update induced by this cost can be accumulated into shared ancestor weights.
% An externally specified feedback-sharing policy determines whether this induced update is admissible for a shared prefix, or whether it must remain local to a branch.
% This distinction matters because treating all induced updates as local wastes reusable statistical signal, while treating all persistent updates as globally shared can expose provider-specific or specialist participation patterns.
% This raises a basic question for workflow-constrained agents: when can the update induced by terminal feedback be reused as a shared hierarchical learning signal, and when should it remain local?

% A natural formal lens for such systems is multi-stage online learning with terminal-only feedback, where a learner selects a full path through a hierarchy and observes only the realized terminal cost~\cite{hou2024distributed}.
% This end-to-end bandit formulation is a robust and conservative starting point: it assumes no more than the cost of the executed path and learns through local bandit estimates along the realized trajectory.
% However, it does not model when the same terminal observation may induce a persistent aggregate update that is reusable by other policies inside a feedback-admissible subtree, while being blocked at structural barriers.
% \textbf{L1: Structural flattening.}
% The realized terminal cost is used only through the sampled trajectory, without distinguishing whether its induced update is also valid for other complete policies sharing a safe prefix or subtree.
% \textbf{L2: Missing barrier-limited sharing.}
% Existing end-to-end bandit formulations do not capture the intermediate regime where the sampled cost can support local bandit learning, while the persistent shared update is allowed to propagate only inside a barrier-free region.

% \emph{When a terminal outcome is admissible for a prefix or subtree under a given feedback-sharing policy, how should a hierarchical learner reuse the induced update without violating structural barriers?}
% Our key insight is that feedback reuse should be governed by structural admissibility.
% Conceptually, terminal feedback can induce aggregate updates inside policy-admissible safe regions, but propagation must stop at ownership, telemetry, provider, or specialist barriers.
% Algorithmically, BarrierShare realizes this principle through an exponential-weight delta.
% When a selected leaf is updated, its accumulated score changes, which induces a change in its exponential weight; this induced weight change defines a \emph{delta} that can be added to admissible ancestor aggregate weights.
% This shared delta differs from the current-round scalar cost used in a local bandit update: over time, accumulated deltas determine the subtree weights used by ancestor decisions, thereby forming a persistent statistical profile of downstream branches.
% When sharing is complete within a subtree, recursive aggregation makes the subtree behave like a single leaf-level adversarial bandit over its leaves, so the regret depends on the number of leaves in that safe subtree rather than accumulating an education penalty across every internal layer.
% When a structural barrier is encountered, the shared delta must stop, and the upstream decision must instead rely on local bandit-style estimation.

% Such barriers are natural in agent workflows with data-silo, provider, contractual, or telemetry restrictions.
% A downstream module may return the current task output and scalar cost, while still prohibiting a globally visible long-run update trace.
% This distinction is motivated by analogous leakage risks in collaborative learning, where shared updates can reveal unintended properties of a participant's data or membership information~\cite{melis2019unintended}, and in gradient-based distributed learning, where gradients can even reconstruct private training examples~\cite{zhu2019deep}.
% It is also relevant to deployed model services, where repeated black-box queries can expose commercially valuable model behavior or enable model extraction~\cite{tramer2016stealing}.
% In sensitive domains such as healthcare, multi-institutional learning is often designed to avoid direct patient-data sharing across silos~\cite{sheller2020federated}.
% By analogy, cumulative shared deltas from a private retrieval specialist, clinical-data module, enterprise knowledge component, or third-party provider may reveal the data domain, provider identity, task distribution, or specialization profile behind that branch.
% We therefore model these interfaces as externally specified feedback-sharing barriers, not as a formal privacy guarantee.
% The relevant regime is neither ``share everything'' nor ``share nothing'': current terminal costs may remain usable through local bandit channels, while persistent shared deltas may propagate upward only until the first structural barrier.
% To address this setting, we propose \textbf{BarrierShare}, a hierarchical online learning algorithm that performs shared delta aggregation inside barrier-free regions and switches to local bandit learning at risky or barrier-limited interfaces.
% Figure~\ref{fig:overview} illustrates this progression from end-to-end bandit feedback to full-share propagation and finally to the barrier-limited regime studied in this paper.

% \textbf{Our contributions are summarized as follows.}
% \textbf{(1) A barrier-limited feedback model.}
% We formalize an intermediate hierarchical feedback regime in which terminal scalar costs remain usable for local bandit learning, but the persistent shared deltas induced by those costs can propagate only within structurally admissible regions.
% This model separates current-round cost feedback from long-run aggregate update traces, and introduces safe regions, structural barriers, and risky depth as the key objects of analysis.
% \textbf{(2) The BarrierShare algorithm.}
% We propose \textbf{BarrierShare}, a single hierarchical online learning algorithm that combines recursive exponential-weight aggregation inside safe subtrees with local bandit updates at barrier-limited risky interfaces.
% A selected shared leaf updates its exponential weight, propagates the induced delta through admissible ancestors, and stops the propagation at the first structural barrier.
% \textbf{(3) Risky-depth regret and validation.}
% We prove regret bounds governed by the risky skeleton depth \(R\), showing how BarrierShare interpolates between fully shared safe aggregation and depth-dependent bandit learning.
% In particular, long safe suffixes reduce the effective depth in the leading bandit term from the full workflow depth \(L\) to the risky depth \(R\), up to safe-subtree aggregation terms.
% We validate these predictions through regret-aligned simulations, including long-safe-suffix and depth-scaling studies, and through fixed-stage agent workflow experiments.

% \section{Related Work}

% BarrierShare is related to three lines of work, distinguished primarily by their information structure.
% First, multi-stage bandits study terminal-only feedback in distributed hierarchies, but use the realized outcome only through local bandit updates.
% Second, feedback-graph, semi-bandit, combinatorial bandit, and partial-monitoring models enrich the current-round observation, whereas BarrierShare keeps terminal-only observations and controls the propagation of the induced update.
% Third, workflow-constrained agents motivate fixed-stage policies and terminal evaluation, but typically optimize execution within a trajectory rather than repeated online selection among workflow policies.
% Table~\ref{tab:feedback_regimes} summarizes these regimes.

% \begin{table}[t]
% \centering
% \small
% \caption{Feedback regimes relevant to BarrierShare. Unlike stronger-feedback models, BarrierShare does not enrich the current-round observation; it restricts the scope over which the terminal-cost-induced update trace may be reused.}
% \label{tab:feedback_regimes}
% \begin{tabular}{lllll}
% \toprule
% Regime & Observation & Reused object & Reuse scope & Example \\
% \midrule
% End-to-end bandit & sampled terminal cost & none & none & \citet{hou2024distributed} \\
% Complete / one-hop & child/action-level feedback & current-round estimates & local node & Hou complete feedback \\
% Full-share endpoint & sampled terminal cost & weight delta & all ancestors & conceptual endpoint \\
% BarrierShare & sampled terminal cost & weight delta & safe subtrees & ours \\
% \bottomrule
% \end{tabular}
% \end{table}

% \subsection{Multi-Stage Bandits and Distributed Online Learning}

% Classical adversarial bandits study a single learner that selects one action and observes only the loss of the selected action, with algorithms such as EXP3 providing no-regret guarantees in the one-stage setting~\citep{auer2002nonstochastic}.
% Multi-stage systems introduce a qualitatively different difficulty: the terminal outcome is determined by a sequence of decisions across stages, and upstream decisions affect not only immediate performance but also whether downstream learners receive enough traffic to learn.
% This creates the education--exploration--exploitation trilemma formalized by \citet{hou2024distributed}.

% The closest theoretical baseline for BarrierShare is Hou's multi-stage bandit model.
% In that formulation, a system is represented as a tree: each internal node selects a child, each leaf produces a terminal cost, and the selected path observes only the realized end-to-end outcome.
% The $\epsilon$-EXP3 algorithm addresses the education problem by mixing uniform exploration with local exponential-weights updates, yielding regret of order \(O(T^{L/(L+1)})\) for an \(L\)-stage system.
% This is the conservative endpoint for our work: terminal feedback is available, but it is used only through local bandit-style updates along the sampled trajectory.

% Hou also studies a stronger complete one-hop feedback setting, where each node can observe richer feedback about all of its actions and normalized-EG-type updates become possible.
% This setting is fundamentally different from BarrierShare: it enriches current-round observations, whereas our model keeps terminal-only scalar feedback and restricts whether the induced persistent update trace can be shared upward.

% The two-stage model of \citet{singla2018learning} is related in that a central forecaster interacts with learning experts and can coordinate their learning by blocking whether the selected expert observes feedback.
% However, this feedback blocking controls the learning state of selected experts in a two-level expert-advice problem, whereas BarrierShare studies arbitrary-depth hierarchies in which terminal costs remain locally usable but the induced persistent update delta may propagate only through structurally admissible ancestors.

% Our setting is also distinct from adversarial MDP and tree-search formulations.
% Although a hierarchy can be viewed as a tree of states, these methods typically assume transition-level state observations, simulator access, value estimates, or centralized coordination.
% Unlike adversarial MDP formulations, our learner does not observe intermediate states or rewards along the trajectory, and the regret comparator is a stationary root-to-leaf workflow policy rather than an adaptive Markov policy.
% Unlike tree-search methods, our tree is a fixed high-level workflow-policy family rather than a low-level action expansion tree, and learning occurs across repeated terminal-cost observations rather than through within-trajectory search.

% \subsection{Feedback Regimes in Hierarchical Learning}

% A broad online learning literature studies how regret changes when the learner receives feedback beyond the loss of the selected action.
% Full-information learning reveals all action losses; semi-bandit and combinatorial bandit models reveal component-level feedback; partial-monitoring and feedback-graph models specify which side observations become available after each action~\citep{cesa2006prediction, mannor2011bandits, audibert2014regret, alon2015online}.
% These settings improve learning by enriching the information available in the current round.

% BarrierShare differs in the object being reused.
% We do not reveal losses of unselected leaves or child-level costs; the only observed scalar remains the terminal cost of the sampled path.
% In a shareable region, this cost induces an exponential-weight delta at the sampled leaf, and this delta updates admissible ancestor aggregates that determine future hierarchical sampling probabilities.
% Thus, the question is not which additional losses are observed in the current round, but how far the induced update trace may propagate.

% Path and combinatorial bandits also consider actions with internal structure, such as paths, subsets, or other combinatorial objects~\citep{cesa2012combinatorial, audibert2014regret}.
% However, semi-bandit variants typically reveal component losses, and bandit variants treat the entire path as a flat arm without providing a mechanism for recursively reusing terminal feedback across safe prefixes.
% BarrierShare lies between these extremes: it never observes component-level losses, but it also avoids flattening the hierarchy whenever the induced update is admissible for a safe subtree.

% As a useful conceptual endpoint, unrestricted full-share recursive aggregation would allow a selected shareable leaf to update all ancestor aggregates through its weight delta.
% BarrierShare retains this mechanism only inside structurally admissible regions; once a barrier is reached, upstream decisions revert to local bandit learning.
% Our motivation is also consistent with observations from collaborative and federated learning that shared updates may expose information about data, participants, or model behavior~\citep{melis2019unintended, zhu2019deep, tramer2016stealing, sheller2020federated}.
% We do not claim a privacy guarantee: barriers are externally specified sharing constraints, and our goal is to characterize their online-learning consequences.

% \subsection{Workflow-Constrained Language Agents}

% Recent language-agent systems increasingly structure execution through tools, search, verifiers, memory, or workflow modules.
% Tool-using and reasoning-acting agents interleave reasoning, actions, and observations when interacting with external environments~\citep{yao2022react, schick2023toolformer}.
% Web-agent and computer-use benchmarks expose long-horizon tasks with tool calls, state tracking, policy constraints, and terminal success criteria~\citep{zhou2023webarena, yao2024tau}.
% These works motivate the view that many agent executions are organized around constrained decision points rather than a single flat action.

% Importantly, the tree in our formulation is not a low-level action tree over clicks, tokens, or simulator states; it is a fixed high-level policy family whose leaves are complete workflow policies or macro-strategies.
% A workflow policy may choose among prompts, tools, specialist modules, API calls, or standard operating procedures at different stages.
% Thus, our abstraction is closer to repeated selection among candidate workflow policies than to online search over low-level environment actions.

% Recent agent work improves trajectory execution through workflow guidance, process reward models and verifiers, value-guided search, memory retrieval, distillation, or computer-use data engines~\citep{liu2018workflow, chae2025webshepherd, mendes2025selftaught, kang2025agentdistillation, zhang2025gmemory, wang2025opencua}.
% These methods change how an agent executes, evaluates, retrieves, or learns within a trajectory.
% Unlike prior agent works that improve the executor, verifier, reward model, search procedure, or memory within a trajectory, BarrierShare studies the learning layer above execution: how to repeatedly select among fixed-stage workflow policies when only terminal costs are observed and only part of the induced update trace can be propagated upward.

% \section{Problem Formulation}

% We study hierarchical online learning on a rooted decision tree under three feedback regimes: the conservative end-to-end bandit regime, an idealized full-share endpoint, and the barrier-limited sharing regime that is the focus of this paper. This section defines the game, the feedback semantics, and the regret objective, without introducing the BarrierShare algorithm itself.

% \paragraph{Notation.}
% Let $\mathcal{T}$ be a rooted tree with root $r$, leaf set $\mathcal{L}$, and child set $C_i$ for each non-leaf node $i$. For any node $i$, let $L(i) \subseteq \mathcal{L}$ denote the set of leaf descendants of $i$. At round $t$, the learner selects a root-to-leaf path $\ell_t \in \mathcal{L}$ by making sequential decisions along the tree. We write $p_t(j \mid i)$ for the probability of selecting child $j \in C_i$ at node $i$, and $\Pi_t(i)$ for the probability of reaching node $i$ under the round-$t$ decision rule. In particular, $\Pi_t(\ell)$ denotes the probability of selecting a leaf $\ell \in \mathcal{L}$. Each leaf $\ell$ incurs a terminal cost $c_t(\ell) \in [0,1]$. For any node $i$, let $Y_\eta[i,T]$ denote the cumulative expected loss of the learner over $T$ rounds starting from $i$, and let $Y_\ast[i,T]$ denote the cumulative loss of the optimal stationary policy starting from $i$.

% \subsection{Multi-Stage Tree Model and Regret Objective}

% We model a multi-stage decision system as a rooted tree $\mathcal{T}$ of depth $L$. Each internal node represents a decision point, and each leaf represents a complete root-to-leaf decision path. At each round $t = 1,\ldots,T$, the learner starts at the root and sequentially selects one child at each internal node until reaching a leaf $\ell_t \in \mathcal{L}$. Equivalently, the round-$t$ decision can be written as a path
% \[
% \ell_t = (a_{1,t}, \ldots, a_{L,t}),
% \]
% where $a_{k,t}$ is the stage-$k$ choice along the selected path.

% The decision rule is specified by local conditional probabilities $p_t(j \mid i)$ at internal nodes. The induced probability of reaching a node is defined recursively by
% \[
% \Pi_t(r) = 1,
% \qquad
% \Pi_t(j) = \Pi_t(i)\, p_t(j \mid i)
% \quad \text{for } j \in C_i.
% \]
% Hence, for a leaf $\ell$, $\Pi_t(\ell)$ is the probability of selecting the entire root-to-leaf path ending at $\ell$.

% After selecting $\ell_t$, the learner incurs and observes the terminal cost $c_t(\ell_t)$. We compare the learner against the best stationary policy that fixes, at every internal node, a single child choice throughout all rounds. Such a policy induces a fixed leaf in every subtree. Define recursively
% \[
% Y_\ast[i,T] =
% \begin{cases}
% \sum_{t=1}^T c_t(i), & i \in \mathcal{L},\\[4pt]
% \min_{j \in C_i} Y_\ast[j,T], & i \notin \mathcal{L}.
% \end{cases}
% \]
% The cumulative expected loss of the learner from node $i$ is
% \[
% Y_\eta[i,T] := \sum_{t=1}^T \mathbb{E}[y_\eta[i,t]],
% \]
% where $y_\eta[i,t]$ is the random loss incurred at round $t$ if the learner starts from $i$ and follows its round-$t$ decision rule thereafter. Our root regret objective is
% \[
% \mathrm{Reg}_T := Y_\eta[r,T] - Y_\ast[r,T].
% \]

% This formulation provides the formal backbone of multi-stage end-to-end learning before any additional feedback structure is introduced.

% \subsection{End-to-End Bandit Regime}

% We first consider the conservative regime in which the learner only observes the scalar terminal cost of the selected path. Formally, after choosing $\ell_t$, the learner observes only
% \[
% c_t(\ell_t),
% \]
% and receives no direct information about the terminal costs of unselected leaves or the intermediate stage-wise contributions along the hierarchy.

% Under this regime, terminal feedback is usable only through bandit-style estimators based on the sampling probability of the selected path or selected local decisions. This model is safe in the sense that it requires no assumption beyond observability of the realized end-to-end outcome. At the same time, it is conservative: even when part of the terminal signal could serve as reusable supervision for an upstream shared structure, the model treats the entire signal as path-local. Representative baselines in this regime include EXP3-type and $\epsilon$-EXP3-type methods, but this subsection defines only the feedback model rather than any particular algorithm.

% \subsection{Full-Share Delta Regime}

% We next define an idealized full-share endpoint in which terminal feedback from a selected leaf can be converted into a sharable delta and propagated upward through the hierarchy.

% Let $\mathcal{L}^{\mathrm{shr}} \subseteq \mathcal{L}$ denote the set of leaves whose terminal feedback is fully shareable with their ancestors. When a selected leaf $\ell_t \in \mathcal{L}^{\mathrm{shr}}$ is observed, its terminal feedback induces a sharable delta, denoted by $\Delta_t(\ell_t)$, which can be transmitted upward to all ancestors that fully observe the corresponding shared subtree.

% The only invariant we require from this regime is the aggregate-visibility property: for any node $i$ that fully observes a shared subtree, its aggregate shared weight equals the sum of the visible descendant shared leaf weights. That is, if $W_t[i]$ denotes the aggregate weight at node $i$, then
% \[
% W_t[i] = \sum_{\ell \in L(i) \cap \mathcal{L}^{\mathrm{shr}}} w_t(\ell),
% \]
% where $w_t(\ell)$ is the shared leaf weight associated with $\ell$. This full-share regime serves as an idealized endpoint that reveals the value of reusable hierarchical feedback, but it does not yet capture structural restrictions on how far such feedback may propagate.

% \subsection{Barrier-Limited Sharing Regime}

% We now define the barrier-limited regime studied in this paper. Each non-leaf node $i$ is equipped with a gate
% \[
% g(i) \in \{\mathrm{share}, \mathrm{unshare}\},
% \]
% which determines whether shared feedback may propagate upward through that node. As before, let $\mathcal{L}^{\mathrm{shr}} \subseteq \mathcal{L}$ denote the set of leaves whose terminal feedback is eligible to initiate shared upload.

% At round $t$, if the selected leaf $\ell_t$ belongs to $\mathcal{L}^{\mathrm{shr}}$, then its terminal feedback may initiate an upward shared upload. However, this propagation continues only until the first ancestor whose gate is $\mathrm{unshare}$. Such a node acts as a structural barrier, beyond which the shared delta cannot propagate.

% A non-leaf node $i$ is called \emph{safe} if the entire relevant subtree beneath $i$ is fully shareable and fully visible at $i$, namely if
% \[
% L(i) \subseteq \mathcal{L}^{\mathrm{shr}}
% \quad \text{and} \quad
% g(u)=\mathrm{share}, \ \forall \text{ non-leaf } u \in \mathrm{subtree}(i).
% \]
% Otherwise, $i$ is called \emph{risky}. Intuitively, a safe node sits above a barrier-free fully shared region, whereas a risky node lies on or above a region where upward propagation may be blocked.

% To quantify the amount of non-shareable structure, define the risky depth recursively by
% \[
% R(i)=
% \begin{cases}
% 0, & \text{if } i \text{ is safe},\\
% 1, & \text{if } i \text{ is risky and } C_i \subseteq \mathcal{L},\\
% 1+\max_{j \in C_i} R(j), & \text{if } i \text{ is risky and } C_i \not\subseteq \mathcal{L}.
% \end{cases}
% \]
% The risky depth of the whole tree is then
% \[
% R := R(r).
% \]
% This quantity captures the maximum number of risky decision points encountered on any root-to-leaf path before entering a barrier-free fully shared suffix.

% \paragraph{Core question.}
% \emph{How can a learner achieve low regret when terminal feedback can be propagated upward only until structural barriers are encountered?}


% \section{BarrierShare}
% \label{sec:barriershare}

% BarrierShare learns how to route probability mass through the tree, how to reuse terminal feedback when sharing is permitted, and how to block propagation when a structural barrier is encountered. The algorithm interpolates between two endpoints: recursive sharing in fully shareable regions and local bandit learning in risky regions.

% \subsection{Method Overview}
% \label{sec:barriershare_overview}

% BarrierShare maintains two coupled mechanisms. The first is \emph{recursive selection}: at each internal node, the learner selects a child either according to aggregated subtree mass or according to a local exploration-exploitation rule. The second is \emph{barrier-limited update}: after observing terminal feedback at the selected leaf, the learner performs local updates on the sampled risky path; if the selected leaf is shareable, it additionally propagates the induced leaf delta upward through the shareable prefix of the path.

% Safe regions reuse feedback through exact aggregation over shared leaves, whereas risky regions only admit local credit assignment from the sampled trajectory. These two mechanisms operate on different state variables but are coupled through the sampled path.

% \subsection{Recursive Selection}
% \label{sec:barriershare_selection}

% Starting from the root, BarrierShare recursively selects one child at each visited internal node until a leaf is reached.

% \paragraph{Safe nodes.}
% For any safe node $u$, child selection is induced by aggregated subtree weights:
% \begin{equation}
%     p_t(v\mid u) = \frac{W_t(v)}{W_t(u)},
%     \qquad v \in \mathsf{ch}(u).
%     \label{eq:safe_selection}
% \end{equation}
% Here $W_t(u)$ is the exact aggregate of shared leaf weights in the shareable subtree rooted at $u$. Because safe subtrees are closed under descendants, the same rule applies recursively to every child subtree below a safe node.

% \paragraph{Risky nodes.}
% If the current node is risky, BarrierShare uses an $\epsilon$-greedy exploration rule:
% \begin{equation}
%     p_t(v\mid u)
%     =
%     \begin{cases}
%     \dfrac{1}{|\mathsf{ch}(u)|}, & \text{with probability } \epsilon_i, \\[8pt]
%     \dfrac{\exp(\eta \theta_t(v\mid u))}{\sum_{v' \in \mathsf{ch}(u)} \exp(\eta \theta_t(v'\mid u))}, & \text{with probability } 1-\epsilon_i,
%     \end{cases}
%     \qquad v \in \mathsf{ch}(u),
%     \label{eq:risky_selection}
% \end{equation}
% where $\epsilon_i \in (0,1)$ controls exploration at node $u$. This rule means that with probability $\epsilon_i$ the learner samples uniformly at random, and with probability $1-\epsilon_i$ it samples according to the exponential-weight distribution.

% The probability of the full sampled path is
% \begin{equation}
%     P_t(\pi_t) \coloneqq \prod_{k=0}^{D_t-1} p_t(u_{k+1}\mid u_k).
%     \label{eq:path_prob}
% \end{equation}
% We use $P_t(\pi_t)$ only to describe the realized trajectory probability; the node-local update below uses the reach probability $\rho_t(u)$ together with the local choice probability $p_t(v\mid u)$.

% \begin{algorithm}[t]
% \caption{\textsc{BarrierShare}}
% \label{alg:barriershare}
% \KwIn{Tree $\mathcal{T}$, shared leaf set $\mathcal{L}_{\mathrm{shr}}$, barrier mask $g$, learning rate $\eta$, exploration rate $\epsilon$}
% \KwOut{Sampled leaf $\ell_t$ and updated states}

% Initialize $\phi_1(\ell)=0$ and $w_1(\ell)=\exp(-\eta \phi_1(\ell))$ for all $\ell\in\mathcal{L}_{\mathrm{shr}}$\;
% Initialize $W_1(u)=\sum_{\ell\in\mathcal{L}(u)\cap \mathcal{L}_{\mathrm{shr}}} w_1(\ell)$ for all safe nodes $u$\;
% Initialize $\theta_1(v\mid u)=0$ for all directed edges $(u,v)$\;

% Set $u \leftarrow r$\;
% \While{$u \notin \mathcal{L}$}{
%     \eIf{$u$ is safe}{
%         sample $v \in \mathsf{ch}(u)$ according to \eqref{eq:safe_selection}\;
%     }{
%         sample $v \in \mathsf{ch}(u)$ according to \eqref{eq:risky_selection}\;
%     }
%     record selected edge $(u,v)$ and set $u \leftarrow v$\;
% }
% Let the sampled leaf be $\ell_t \leftarrow u$ and observe terminal cost $c_t$\;

% \ForEach{risky node $u$ on $\pi_t$ with sampled child $v$}{
%     compute
%     \[
%         \widehat{c}_t(u,v)
%         \coloneqq
%         \frac{c_t \cdot \mathbf{1}\{(u,v)\in\pi_t\}}{\rho_t(u)\,p_t(v\mid u)}\;,
%     \]
%     and update
%     \[
%         \theta_{t+1}(v\mid u) = \theta_t(v\mid u) - \eta \widehat{c}_t(u,v)\;.
%     \]
% }

% \eIf{$\ell_t \in \mathcal{L}_{\mathrm{shr}}$}{
%     update $\phi_{t+1}(\ell_t)=\phi_t(\ell_t)+c_t$ and recompute $w_{t+1}(\ell_t)=\exp(-\eta\phi_{t+1}(\ell_t))$\;
%     compute the induced leaf delta $\Delta_t(\ell_t)=w_{t+1}(\ell_t)-w_t(\ell_t)$\;
%     additionally propagate $\Delta_t(\ell_t)$ upward along ancestors until the first node with $g(u)=\mathrm{unshare}$, and stop there\;
% }{
%     do not generate any reusable ancestor update\;
% }
% \end{algorithm}

% \subsection{Risky Local Update}
% \label{sec:barriershare_risky_update}

% The risky update is a node-local bandit update. Let $u$ be a risky node on the sampled path and let $v$ be the chosen child of $u$. Since the learner reaches $u$ with probability $\rho_t(u)$ and then chooses $v$ with probability $p_t(v\mid u)$, we use the importance-weighted estimator
% \begin{equation}
%     \widehat{c}_t(u,v)
%     \coloneqq
%     \frac{c_t \cdot \mathbf{1}\{(u,v)\in\pi_t\}}{\rho_t(u)\,p_t(v\mid u)}.
%     \label{eq:iw_estimator}
% \end{equation}
% The local score update is
% \begin{equation}
%     \theta_{t+1}(v\mid u)
%     =
%     \theta_t(v\mid u) - \eta \widehat{c}_t(u,v).
%     \label{eq:risky_local_update}
% \end{equation}

% This update only affects the visited risky edge. It does not modify subtree weights above the barrier, and it does not assume that the terminal cost is reusable outside the sampled local interface.

% \subsection{Shared Leaf Update and Delta Propagation}
% \label{sec:barriershare_shared_update}

% If the selected leaf $\ell_t$ is shareable, then the terminal feedback is also eligible for reusable aggregation. We update the leaf score by
% \begin{equation}
%     \phi_{t+1}(\ell_t) = \phi_t(\ell_t) + c_t,
%     \label{eq:leaf_score_update}
% \end{equation}
% which induces the updated leaf weight
% \begin{equation}
%     w_{t+1}(\ell_t) = \exp\!\bigl(-\eta \phi_{t+1}(\ell_t)\bigr).
%     \label{eq:leaf_weight_update}
% \end{equation}
% We then define the induced leaf delta as
% \begin{equation}
%     \Delta_t(\ell_t) \coloneqq w_{t+1}(\ell_t) - w_t(\ell_t).
%     \label{eq:delta_def}
% \end{equation}

% For every ancestor $u$ of $\ell_t$ such that the path from $u$ to $\ell_t$ contains only share-labeled nodes, we propagate $\Delta_t(\ell_t)$ upward and refresh the corresponding subtree weights. The propagation stops at the first barrier node, and no ancestor above that node receives the reusable update.

% This rule keeps multiplicative leaf updates and additive subtree aggregation consistent: leaf scores are updated through the exponential map, while subtree weights are maintained as exact sums of shared leaf weights. In particular, for every ancestor that receives the propagated delta, the updated subtree weight remains equal to the sum of the updated shared descendant leaf weights.

% If $\ell_t \notin \mathcal{L}_{\mathrm{shr}}$, then no reusable ancestor update is generated. The terminal cost still contributes to the risky local update on the sampled path, but it is not propagated into any subtree weight.

% \subsection{Endpoint Reductions}
% \label{sec:barriershare_reductions}

% \paragraph{All-share regime.}
% If every node is shareable and every leaf is shared, then every internal node is safe. In this case, BarrierShare reduces to a fully recursive sharing mechanism in which all routing probabilities are induced by subtree weights and every terminal feedback update propagates to the root.

% \paragraph{All-unshare regime.}
% If no reusable sharing is available, then every non-root node is risky. In this case, BarrierShare reduces to a purely local hierarchical bandit learner that updates only the sampled risky edges along the path.

% \paragraph{Intermediate regime.}
% In the general case, safe regions reuse shared feedback through exact subtree aggregation, while risky regions remain local and are updated only through bandit feedback. This is the regime of primary interest, because it interpolates between full sharing and no sharing in a structurally interpretable way.

% \subsection{Application to Hierarchical Agent Workflows}
% \label{sec:barriershare_agent_instantiation}

% The same formalism naturally instantiates hierarchical agent workflows. In that setting, each internal node corresponds to a stage-level routing decision and each leaf corresponds to a complete terminal macro-policy. The learner selects which leaf policy to execute, while the executor carries out the leaf-level trajectory and returns a single terminal cost after completion.

% This mapping separates routing from execution. BarrierShare learns how to route probability mass through the tree, but it does not require step-wise supervision inside each leaf. Therefore, the same algorithm can be used for synthetic tree simulations and for real hierarchical agent pipelines, provided that the terminal cost is observed only after the selected leaf terminates.


% \section{Experiments}

% Our experiments are designed to examine whether barrier-limited feedback sharing is empirically useful, and whether the observed behavior is consistent with the structural claim of the paper. The claim under test is that, in a hierarchical decision process, feedback can be shared upward only within structurally valid regions, while the effective learning difficulty is governed primarily by the risky portion of the tree rather than by total depth alone.

% We evaluate this claim in two complementary settings: a fixed-stage LLM workflow and synthetic tree simulation. The workflow setting tests whether the mechanism remains meaningful in a realistic staged decision process, while the simulation setting isolates the structural variables and allows direct testing of the scaling behavior predicted by the theory.

% We organize the experiments around three questions:
% \textbf{(Q1)} In a fixed workflow with the same budget and evaluator, does BarrierShare behave differently from endpoint sharing regimes and standard bandit baselines?
% \textbf{(Q2)} In simulation, do learning curves vary more strongly with risky depth and barrier placement than with total depth alone?
% \textbf{(Q3)} Does the method remain stable under barrier misspecification, perturbation, and changes in branching factor?

% \subsection{Evaluation Principles}

% We evaluate all methods under the same tree structure, the same horizon, the same random-seed schedule, and the same terminal evaluation protocol within each comparison group. For workflow experiments, we also keep the same executor and the same prompt template family across methods. For simulation experiments, all methods are run on the same generated tree instances.

% We report mean and standard deviation over multiple seeds, and we include learning curves whenever possible rather than only final values. For simulation, the primary metric is cumulative regret $\mathrm{Reg}_T$ and normalized regret $\mathrm{Reg}_T/T$. For the workflow setting, the primary metric is success rate, with average total cost as the main efficiency metric. Token usage, number of calls, latency, and API cost are reported as secondary diagnostics.

% Because the core claim of the paper concerns structure, we also report intermediate quantities that expose the mechanism directly. These include barrier stop depth, propagation depth, shared-update fraction, and the amount of update mass reaching ancestor nodes. These diagnostics are used as mechanism checks, not as replacements for the main metrics.

% \paragraph{Fairness note.}
% BarrierShare is compared only against methods that operate under the same observable tree and the same terminal feedback protocol. If a baseline does not use barrier structure, that is made explicit. If a baseline is an endpoint regime, it is treated as a structural reference point rather than as a fully realistic deployment policy.

% \subsection{Workflow Instantiation}

% To connect the hierarchical model to a practical staged process, we instantiate BarrierShare in a fixed-stage workflow. Each stage corresponds to one decision layer in the tree, and each complete root-to-leaf path corresponds to one workflow policy. In the telecom MMS setting used in our experiments, the workflow contains five stages:
% (1) user grounding,
% (2) customer-line resolution,
% (3) observed feature extraction,
% (4) blocker adjudication and repair planning, and
% (5) terminal execution decision.

% This mapping is intended to make the structural assumptions observable. Some feedback in this workflow is reusable across multiple ancestors because it concerns generic routing, common prefixes, or broadly valid diagnosis. Other feedback is branch-local because it depends on provider-specific state, specialist behavior, or execution details. The barrier defines where upward propagation stops. It is therefore a modeling constraint, not a privacy claim and not an oracle guarantee.

% To reduce ambiguity, we define the barrier map before evaluation and keep it fixed during a run. The map is constructed from the workflow specification rather than from terminal outcomes.
% \textcolor{gray}{[TODO: add the precise barrier construction rule, e.g. how shared vs. unshared regions are specified and how the barrier map is derived from the workflow.]}

% \subsection{Methods Compared}

% We compare BarrierShare with the following baselines:

% \begin{itemize}
%     \item \texttt{full-share}: an endpoint regime in which feedback is propagated upward whenever the structure allows it, without barrier restriction.
%     \item \texttt{full-unshare}: an endpoint regime in which feedback remains local and no upward sharing is performed.
%     \item EXP3: a standard bandit baseline on the same tree structure.
%     \item $\epsilon$-EXP3: a local exploration baseline with explicit randomization.
%     \item Random: a sanity-check baseline.
%     \item Naive: a simple heuristic baseline with no structured propagation rule.
% \end{itemize}

% We treat \texttt{full-share} and \texttt{full-unshare} as structural endpoints, not as identical algorithmic peers of BarrierShare. They are included because they isolate the two extremes of the feedback-sharing assumption. EXP3 and $\epsilon$-EXP3 provide standard bandit references, while Random and Naive are included only as weak sanity checks.

% \textcolor{gray}{[TODO: add one sentence clarifying the exact information available to each baseline, to make the comparison protocol fully explicit.]}

% \subsection{Main Workflow Results}

% \begin{table}[t]
%     \centering
%     \small
%     \caption{Workflow results. We report mean $\pm$ standard deviation over multiple seeds. Lower cost is better and higher success is better.}
%     \label{tab:llm_main}
%     \begin{tabular}{llcccc}
%         \toprule
%         Setting & Method & Success $\uparrow$ & Cost $\downarrow$ & LLM Calls $\downarrow$ & Tokens $\downarrow$ \\
%         \midrule
%         Fixed-stage workflow & BarrierShare & -- & -- & -- & -- \\
%         Fixed-stage workflow & \texttt{full-share} & -- & -- & -- & -- \\
%         Fixed-stage workflow & \texttt{full-unshare} & -- & -- & -- & -- \\
%         Fixed-stage workflow & EXP3 & -- & -- & -- & -- \\
%         Fixed-stage workflow & $\epsilon$-EXP3 & -- & -- & -- & -- \\
%         Fixed-stage workflow & Random & -- & -- & -- & -- \\
%         Fixed-stage workflow & Naive & -- & -- & -- & -- \\
%         \bottomrule
%     \end{tabular}
% \end{table}

% The workflow results are interpreted conservatively.
% \textcolor{gray}{[TODO: after results are available, add a neutral description of the main comparison pattern here. Avoid strong claims until the numbers are filled in.]}

% We therefore examine whether the final outcome is accompanied by reduced unnecessary propagation, lower update mass on barrier-violating ancestors, and a more concentrated use of shared prefixes.
% \textcolor{gray}{[TODO: add the result-dependent interpretation after the table is filled.]}

% \subsection{Workflow Mechanism Checks}

% To separate the effect of propagation from the effect of generic task performance, we log several intermediate variables during workflow execution.

% First, we measure \emph{barrier stop depth}, which indicates where shared updates cease to propagate. Second, we measure \emph{shared-update fraction}, which tracks how often updates are assigned to shared ancestors rather than branch-local nodes. Third, we measure \emph{propagation depth}, which records how far a terminal signal travels upward before it is stopped by a barrier. Fourth, we measure \emph{risky-path frequency}, which indicates how often the execution enters a region where feedback is not safely reusable.

% \textcolor{gray}{[TODO: specify exactly how each diagnostic is computed and aggregated across runs.]}

% \subsection{Simulation: Risky Depth and Safe Suffix}

% We use synthetic simulation to test the structural prediction under direct control of the tree variables. This setting is where the main theoretical claim can be examined most cleanly. The goal is to determine whether the observed regret depends primarily on the risky portion of the tree and on the location of the barrier, rather than on total depth alone.

% \paragraph{Tree construction.}
% We generate rooted trees with controlled depth, branching factor, and share pattern.
% \textcolor{gray}{[TODO: add the exact tree-generation procedure, including how share patterns are sampled or constructed.]}
% We vary the risky skeleton depth $R$ while holding total depth $L$ fixed, and we also vary $L$ while holding the risky structure fixed. This allows us to compare the cost of longer hierarchies in general and the cost of the non-shared portion specifically.

% \paragraph{Primary simulation table.}
% \begin{table}[t]
%     \centering
%     \small
%     \caption{Simulation results under controlled tree structures. We report mean $\pm$ standard deviation over multiple seeds.}
%     \label{tab:sim_main}
%     \begin{tabular}{llccc}
%         \toprule
%         Tree & Method & $\mathrm{Reg}_T \downarrow$ & $\mathrm{Reg}_T/T \downarrow$ & Learning curve fit \\
%         \midrule
%         all-share & BarrierShare & -- & -- & -- \\
%         partial-share & BarrierShare & -- & -- & -- \\
%         all-unshare & BarrierShare & -- & -- & -- \\
%         partial-share & \texttt{full-share} & -- & -- & -- \\
%         partial-share & \texttt{full-unshare} & -- & -- & -- \\
%         partial-share & EXP3 & -- & -- & -- \\
%         partial-share & $\epsilon$-EXP3 & -- & -- & -- \\
%         partial-share & Random & -- & -- & -- \\
%         partial-share & Naive & -- & -- & -- \\
%         \bottomrule
%     \end{tabular}
% \end{table}

% \paragraph{Scaling analysis.}
% To make the structural claim testable, we analyze how regret changes with $R$, $L$, and branching factor. We report learning curves and fitted scaling trends, rather than only endpoint values.
% \textcolor{gray}{[TODO: after full results are available, add the exact scaling statement supported by the plots.]}

% \paragraph{Endpoint sanity checks.}
% We also include endpoint settings as sanity checks. In the all-share case, the full tree should behave like a fully shared region. In the all-unshare case, the method should reduce to local learning without upward reuse.

% \subsection{Barrier Perturbation and Misspecification}

% We next test whether the method is sensitive to errors in the barrier map. This is a key robustness experiment because barrier-limited propagation relies on a structural assumption about which nodes should be allowed to receive shared updates.

% We consider three perturbation settings.

% \paragraph{False positive barriers.}
% A node that should allow upward sharing is incorrectly marked as a barrier.

% \paragraph{False negative barriers.}
% A node that should stop propagation is incorrectly treated as shareable.

% \paragraph{Granularity mismatch.}
% The barrier map is specified at the wrong structural level, for example at the stage level instead of the branch level, or at the branch level instead of the subtree level.

% \textcolor{gray}{[TODO: add the perturbation rates or corruption levels once the final experimental design is fixed.]}

% \subsection{Candidate-Set Scaling}

% We treat candidate-set scaling as a secondary robustness analysis. Here we vary the branching factor while keeping the barrier structure fixed. This tests whether BarrierShare remains useful when each stage has many candidate actions.

% \textcolor{gray}{[TODO: specify the branching-factor settings and whether this analysis stays in the main paper or moves to the appendix.]}

% \subsection{Efficiency and Resource Use}

% Because the workflow setting is resource-sensitive, we report efficiency metrics alongside success. The main efficiency metric is average total cost. We also report cost per success and tokens per success when these quantities are well-defined. For completeness, we include the number of LLM calls, input tokens, output tokens, and average latency.

% \textcolor{gray}{[TODO: once the workflow numbers are available, add the final efficiency interpretation here.]}

% \subsection{Additional Robustness Checks}

% We include two additional robustness analyses.

% \paragraph{Specialist mass.}
% We vary the mass assigned to private or specialist branches. This test examines whether the method remains useful when a larger fraction of the tree lies in regions where upward sharing is not valid.

% \paragraph{Noise and execution instability.}
% We perturb terminal outcomes with bounded noise to test whether the method remains stable under imperfect execution.

% \textcolor{gray}{[TODO: decide whether these analyses stay in the main paper or are moved to the appendix.]}

% \subsection{What the Experiments Are Meant to Establish}

% The experiments are designed to compare three possible interpretations of the results:
% \begin{itemize}
%     \item BarrierShare behaves like a generic hybrid heuristic and its structure is not essential.
%     \item BarrierShare is useful only in a narrow subset of favorable settings.
%     \item BarrierShare exhibits behavior that is consistent with barrier-limited sharing, with performance depending more strongly on risky depth and barrier placement than on total depth alone.
% \end{itemize}

% \textcolor{gray}{[TODO: after results are filled in, add a neutral concluding paragraph summarizing which interpretation is most consistent with the data.]}


% \section{Conclusion}

% This paper studies hierarchical online learning for multi-stage agent workflows under structured feedback sharing. We begin from the conservative end-to-end bandit view, show that full-share delta propagation reveals the value of reusable upward feedback, and then identify structural barriers as the key reason why real systems lie between full sharing and purely path-local learning. In response, we propose BarrierShare, a unified intermediate algorithm that combines recursive shared aggregation in barrier-free regions with local bandit learning at risky decision points. Our theoretical results show that the resulting learning difficulty is governed by the risky skeleton depth rather than only by the total hierarchy depth, and our simulation and fixed-stage agent experiments provide consistent evidence for this structural picture. A current limitation is that our formulation is tailored to fixed-stage or fixed-for-the-experiment hierarchical workflows, where barriers are modeled as structural feedback-sharing constraints rather than formal privacy guarantees.

\bibliographystyle{unsrtnat}
\bibliography{main2_refs}
\section{Notation and Additional Definitions}

\towrite{Use this appendix to reduce notation load in the main paper and make the proofs easier to read. Keep it concise but systematic. This section should gather the symbols that are currently spread across the problem setup, method, and theory.}

\subsection{Global Symbol Table}

\towrite{Provide a compact table for the main symbols: tree objects $(\mathcal{T}, r, \mathcal{L}, C_i, L(i))$; path and probability symbols $(\ell_t, \Pi_t(\cdot), p_t(j \mid i))$; loss and comparator symbols $(c_t(\ell), Y_\eta[i,T], Y_\ast[i,T], C_i^\ast, d_i)$; and algorithm-state symbols $(\theta_\ell, W[u], \Delta_t)$.}

\subsection{Barrier and Sharing Definitions}

\towrite{Collect the paper-specific definitions in one place: shared leaf, unshared leaf, gate, barrier, upload-reachable edge, safe node, risky node, risky height, risky skeleton depth $R$, and safe suffix. This should be the canonical reference for the appendix proofs.}

\section{Complete Algorithms and Pseudocode}

\towrite{This appendix should contain the full algorithmic descriptions that are too bulky for the main paper. Keep the main text focused on the conceptual modules, and use this appendix for exact procedural definitions.}

\subsection{Full-Share Algorithm}

\towrite{Provide the complete recursive aggregation pseudocode for the full-share endpoint, including the leaf update, aggregate update, and the invariant maintained by internal nodes.}

\subsection{Risky-PS and Intermediate Variants}

\towrite{Include pseudocode for the risky or truncated partial-share variant if you want to show how the endpoint and intermediate algorithms differ. This is the right place to expose the relationship among the variants without cluttering the main text.}

\subsection{BarrierShare Algorithm}

\towrite{Provide the full BarrierShare pseudocode: path selection, local bandit updates on risky prefixes, shared delta propagation, and barrier stopping. This should be the appendix version that a reader can directly implement.}

\subsection{Implementation Notes}

\towrite{Add short implementation notes for choices that matter in code but are too detailed for the main text, such as initialization, probability floors, shared-upload bookkeeping, or how the theorem-facing notation maps to actual state variables.}

\section{Full Theoretical Proofs}

\towrite{This is one of the core appendices. The main paper should keep only theorem statements, key lemmas, and proof intuition; this appendix should contain the complete derivations.}

\subsection{Problem-Dependent Definitions}

\towrite{Restate the definitions needed by the proofs, including safe subtrees, risky nodes, risky height, risky skeleton depth $R$, and barrier-free fully shared suffixes. This subsection should make the later proofs self-contained.}

\subsection{Safe-Subtree Equivalence Proof}

\towrite{Give the full proof that recursive aggregation on a fully safe subtree is equivalent to leaf-level exponential weights or an Exp3-style forecaster over the subtree leaves. The telescoping probability argument belongs here in full detail.}

\subsection{Local Risky-Node Regret Bound}

\towrite{Provide the complete risky single-layer regret derivation, including the branch-conditioned bandit update and any truncated or barrier-specific adjustments. If the main text only presents the final bound, all intermediate steps should be placed here.}

\subsection{Root Regret Proof}

\towrite{Provide the complete proof of the main theorem. This subsection should show how the recursive excess term $\Delta(i)$ is constructed, how the $R$-adjusted quantity enters the bound, and how the proof recovers the all-share and all-unshare endpoints.}

\subsection{Structural Corollaries}

\towrite{Expand the theorem into the corollaries used by the experiments: long safe suffix, depth effect, and the behavior when $R<L$. If you have additional structural discussion about coefficients or branch factors, place it here.}

\subsection{Optional Lower-Bound or Endpoint Discussion}

\towrite{If you decide to include lower-bound comparisons, naive-fix failures, or a more technical discussion of how BarrierShare relates to Hou-style worst-case behavior, put that material here instead of in the main text.}

\section{Experimental Setup}

\towrite{This appendix should make the experiments reproducible. Because your paper mixes LLM-backed workflow experiments and simulation, the setup appendix needs clear subsections rather than one long block of text.}

\subsection{Common Setup}

\towrite{Describe shared configuration choices: seeds, horizon or number of rounds, averaging procedure, regret aggregation, mean and standard deviation reporting, and any common evaluation conventions used across both experiment families.}

\subsection{Fixed-Stage LLM Experiment Setup}

\towrite{Specify how the fixed-stage workflow is instantiated: the stages, candidate agents or policies, terminal cost components, workflow-specific barrier semantics, and how token, cost, and latency statistics are collected. This is where you can document executor details without overloading the main paper.}

\subsection{Simulation Setup}

\towrite{Describe tree-generation rules, branching factors, depth, share-ratio construction, long-safe-suffix construction, specialist-mass construction, and noise injection. This subsection should make the simulation regime fully unambiguous.}

\subsection{Baselines}

\towrite{Describe BarrierShare, EXP3, epsilon-EXP3, Random, Naive, and any endpoint configurations such as full-share or full-unshare. This is the right place for implementation-level baseline details that do not belong in the main text.}

\section{Additional Simulation Results}

\towrite{Use this appendix for complete simulation reporting. The main paper should keep only the most decision-relevant tables and figures; this appendix should contain the full numerical picture.}

\subsection{Full Simulation Tables}

\towrite{Provide the full simulation result tables for all methods, share ratios, seeds, and summary statistics. This subsection should include the complete version of the main simulation table from the paper.}

\subsection{Long Safe Suffix Results}

\towrite{Provide the detailed results for varying risky depth $R$ under fixed total depth $L$, including full tables, curves, or error bars if available. This subsection should make the theorem-to-experiment connection explicit.}

\subsection{Depth Scaling Results}

\towrite{Provide the complete results for depth sweeps, such as $L=5,15,25$, including all methods and all share ratios that are too bulky for the main paper.}

\subsection{Specialist-Mass and Noise Results}

\towrite{If these experiments are run, place them here as mechanism and robustness supplements. They are useful, but they usually do not need to compete for space in the main paper.}

\subsection{Variance and Error Bars}

\towrite{Provide per-seed summaries, confidence intervals, or additional variance plots if they help establish stability. This is especially useful when the simulation evidence is part of the main theorem support.}

\section{Additional LLM Results}

\towrite{Use this appendix for all workflow-backed experimental material that is too detailed for the main paper. The main text should show the headline LLM table and the core ablations; this appendix should carry the rest.}

\subsection{Full Main-Result Tables}

\towrite{Provide the complete version of the main LLM result table, including all methods, share ratios, and metrics that may have been compressed in the main paper.}

\subsection{Interaction Mechanism Ablations}

\towrite{Provide the full ablation tables for algorithm-direct, theta-guided-agent, and agent-only settings, together with any additional breakdowns that do not fit in the main text.}

\subsection{Candidate-Set Scaling Results}

\towrite{Provide the full candidate-scaling tables for 5, 15, and 25 candidates per stage, including all methods and any variance statistics.}

\subsection{Efficiency Analysis}

\towrite{This subsection should contain the full efficiency-oriented reporting: input tokens, output tokens, total tokens, estimated LLM cost, average latency, p95 latency, tokens per success, and cost per success.}

\subsection{Optional Case Studies}

\towrite{If you include qualitative examples, use this subsection for success and failure cases showing why a barrier helped or failed to help. Keep the main paper quantitative and reserve narrative examples for here.}

\section{Fixed-Stage Agent Workflow Instantiation Details}

\towrite{The main paper should only include the minimal mapping needed to understand the experiments. This appendix should contain the full workflow instantiation details needed for careful reading or reproduction.}

\subsection{Tree Schema and Stage Semantics}

\towrite{Describe the fixed-stage schema in full detail: the meaning of each stage, what each node or candidate represents, and how leaves correspond to complete workflow or macro-policies.}

\subsection{Terminal Cost and Barrier Semantics}

\towrite{Document the full terminal-cost definition and explain how shared, unshared, and barrier-limited semantics are interpreted in the workflow environment. This is the right place to spell out workflow-specific nuances that would be too heavy for the main paper.}

\subsection{Additional Environment Details}

\towrite{Add any remaining fixed-tree environment details that a reviewer may want to inspect, such as task-family construction, stage-specific assumptions, or how workflow constraints are encoded.}

\section{Limitations and Future Directions}

\towrite{Keep a short but explicit appendix section on limitations. Even if the main paper no longer has a standalone discussion section, the appendix should still document the paper's scope and boundaries.}

\subsection{Current Limitations}

\towrite{State the main limitations clearly: BarrierShare is designed for fixed-stage or fixed-for-the-experiment hierarchical workflows; the barrier is a structural feedback model rather than a formal privacy guarantee; and the comparator is a best fixed path rather than a dynamic comparator.}

\subsection{Future Directions}

\towrite{Outline the most natural extensions: dynamic trees, planner-tree co-design, learned barriers, richer feedback models, or extending beyond fixed-stage workflows toward more open-ended planning systems.}

\end{document}
